\documentclass[12pt]{amsart}
\usepackage{luatex85}
\usepackage{luatexja-fontspec}
\IfFontExistsTF{Noto Serif CJK JP}
  {\setmainjfont{Noto Serif CJK JP}}
  {\setmainjfont{HaranoAjiMincho-Regular}}
\ltjsetparameter{jacharrange={-2}}


% ============================================================
% Basic spacing
% ============================================================
\setlength{\lineskip}{0pt}

% ============================================================
% Packages for mathematics
% ============================================================
\usepackage{amsmath}
\usepackage{mathtools}
\usepackage{amssymb}
\usepackage{amsthm}
\usepackage{amscd}
\usepackage{mathrsfs}
% Unused package unavailable here: dsfont
% Unused package unavailable here: bbm

% ============================================================
% Graphics
% Use the following line for pLaTeX/upLaTeX + dvipdfmx.
% If you compile with pdfLaTeX or LuaLaTeX, remove the option [dvipdfmx].
% For PNG/JPG files with dvipdfmx, run extractbb if LaTeX cannot find the bounding box.
% ============================================================
%\usepackage[dvipdfmx]{graphicx}
\usepackage{graphicx} % Use this instead for pdfLaTeX or LuaLaTeX.

% ============================================================
% Packages for colors and text decorations
% ============================================================
\usepackage[dvipsnames]{xcolor}
% Unused package unavailable here: ulem
\usepackage{comment}

% ============================================================
% Packages for tables, lists, and diagrams
% ============================================================
\usepackage{multirow}
\usepackage{bigdelim}
\usepackage{enumerate}
\usepackage{enumitem}
\usepackage[all]{xy}
\usepackage{tikz}





% ============================================================
% tcolorbox settings
% Use as: \begin{tcolorbox}[mybox] ... \end{tcolorbox}
% ============================================================
\usepackage[most]{tcolorbox}
\tcbuselibrary{breakable, skins, theorems}
\tcbset{
  mybox/.style={
    colback=white,
    colframe=green!35!black,
    fonttitle=\bfseries,
    breakable=true
  }
}



% ============================================================
% Draft tools
% Comment out showkeys before submitting to a journal or arXiv.
% ============================================================
\usepackage[final]{showkeys} % Use this when you want to hide all labels.
%\usepackage{showkeys} % Use this when you want to show labels.
\renewcommand*{\showkeyslabelformat}[1]{%
  \begingroup
  \fboxsep=2pt%
  \fboxrule=0.3pt%
  \fbox{%
    \parbox{1.8cm}{%
      \raggedright
      \normalfont\tiny\sffamily
      \strut #1\strut
    }%
  }%
  \endgroup
}
%

% ============================================================
% This setting makes every enumerate environment use labels of the form (1), (2), (3), ...
% ============================================================

% Duplicate enumitem load removed.
\setlist[enumerate]{label=$(\arabic*)$}

% ============================================================
% Hyperlinks
% Load hyperref near the end of the package list.
% Use the first line for pLaTeX/upLaTeX + dvipdfmx.
% If you compile with pdfLaTeX or LuaLaTeX, use the second line instead.
% ============================================================
%\usepackage[dvipdfmx,colorlinks,linkcolor=red,anchorcolor=blue,citecolor=blue]{hyperref}
\usepackage[colorlinks,linkcolor=red,anchorcolor=blue,citecolor=blue]{hyperref}



% ============================================================
% Page layout
% ============================================================
\setlength{\topmargin}{-0.25cm}
\setlength{\textwidth}{15cm}
\setlength{\textheight}{22.6cm}
\setlength{\headheight}{1em}
\setlength{\headsep}{0.5cm}
\setlength{\oddsidemargin}{0.40cm}
\setlength{\evensidemargin}{0.40cm}
\setlength{\baselineskip}{15pt}
\setlength{\footskip}{32pt}

% ============================================================
% Theorem environments
% ============================================================
\newtheorem{thm}{Theorem}[section]
\newtheorem{theo}[thm]{Theorem}
\newtheorem{cor}[thm]{Corollary}
\newtheorem{prop}[thm]{Proposition}
\newtheorem{conj}[thm]{Conjecture}
\newtheorem*{mainthm}{Theorem}

\newtheorem{deflem}[thm]{Definition-Lemma}
\newtheorem{lem}[thm]{Lemma}

\theoremstyle{definition}
\newtheorem{defn}[thm]{Definition}
\newtheorem{propdefn}[thm]{Proposition-Definition}
\newtheorem{lemdefn}[thm]{Lemma-Definition}
\newtheorem{thmdefn}[thm]{Theorem-Definition}
\newtheorem{eg}[thm]{Example}
\newtheorem{ex}[thm]{Example}
\newtheorem{exa}[thm]{Example}
\newtheorem{ques}[thm]{Question}
\newtheorem{remin}[thm]{Reminder}

\theoremstyle{remark}
\newtheorem{rem}[thm]{Remark}
\newtheorem{setup}[thm]{Setup}
\newtheorem{obs}[thm]{Observation}
\newtheorem{notation}[thm]{Notation}
\newtheorem{claim}[thm]{Claim}
\newtheorem{assup}[thm]{Assumption}
\newtheorem{cln}[thm]{Claim}

% Independent theorem-like environments.
\newtheorem{cl}{Claim}
\newtheorem{step}{Step}
\newtheorem{case}{Case}

% Unnumbered theorem-like environments.
\newtheorem*{clproof}{Proof of Claim}
\newtheorem*{ack}{Acknowledgements}

% Number equations by section.
\numberwithin{equation}{section}

% ============================================================
% Mathematical operators
% ============================================================
\newcommand{\rk}{\operatorname{rk}}
\newcommand{\rank}{\operatorname{rank}}
\newcommand{\degree}{\operatorname{deg}}
\newcommand{\codim}{\operatorname{codim}}
\newcommand{\nd}{\operatorname{nd}}

\newcommand{\Supp}{\operatorname{Supp}}
\newcommand{\supp}{\operatorname{Supp}}
\newcommand{\Rad}{\operatorname{Rad}}
\newcommand{\Exc}{\operatorname{Exc}}
\newcommand{\Div}{\operatorname{div}}

\newcommand{\End}{\operatorname{End}}
\newcommand{\eend}{\operatorname{End}}
\newcommand{\Coker}{\operatorname{Coker}}
\newcommand{\GL}{\operatorname{GL}}

\newcommand{\Hom}{\mathscr{H}\!\mathit{om}}
\newcommand{\SheafHom}[1]{\mathscr{H}\!\mathit{om}_{#1}}
\newcommand{\PGL}{\mathbb{P}\GL(r,\C)}

\DeclareMathOperator{\Ric}{Ric}
\DeclareMathOperator{\Vol}{Vol}
\DeclareMathOperator{\tr}{tr}
\DeclareMathOperator{\Id}{Id}
\DeclareMathOperator{\res}{res}
\DeclareMathOperator{\length}{length}
\DeclareMathOperator{\Homop}{Hom}
\DeclareMathOperator{\Diff}{Diff}
\DeclareMathOperator{\Lie}{Lie}
\DeclareMathOperator{\Autop}{Aut}
\DeclareMathOperator{\Kerop}{Ker}
\DeclareMathOperator{\Imageop}{Im}


%These are added by TOMOYUKI 

\newcommand{\MY}{\operatorname{MY}}
\newcommand{\abs}[1]{\left\lvert#1\right\rvert}
\newcommand{\norm}[1]{\left\|#1\right\|} 
\newcommand{\Rm}{\operatorname{Rm}}
\newcommand{\Cal}{\operatorname{Cal}}
\newcommand{\NA}{\operatorname{NA}}

\newcommand{\cA}{\mathcal{A}}
\newcommand{\cB}{\mathcal{B}}
\newcommand{\cC}{\mathcal{C}}
\newcommand{\cD}{\mathcal{D}}
\newcommand{\cE}{\mathcal{E}}
\newcommand{\cF}{\mathcal{F}}
\newcommand{\cG}{\mathcal{G}}
\newcommand{\cH}{\mathcal{H}}
\newcommand{\cI}{\mathcal{I}}
\newcommand{\cJ}{\mathcal{J}}
\newcommand{\cK}{\mathcal{K}}
\newcommand{\cL}{\mathcal{L}}
\newcommand{\cM}{\mathcal{M}}
\newcommand{\cN}{\mathcal{N}}
\newcommand{\cO}{\mathcal{O}}
\newcommand{\cP}{\mathcal{P}}
\newcommand{\cQ}{\mathcal{Q}}
\newcommand{\cR}{\mathcal{R}}
\newcommand{\cS}{\mathcal{S}}
\newcommand{\cT}{\mathcal{T}}
\newcommand{\cU}{\mathcal{U}}
\newcommand{\cW}{\mathcal{W}}
\newcommand{\cX}{\mathcal{X}}
\newcommand{\cY}{\mathcal{Y}}
\newcommand{\cZ}{\mathcal{Z}}

% ============================================================
% Standard maps and geometric notation
% ============================================================
\DeclareMathOperator{\pr}{pr}
\DeclareMathOperator{\id}{id}
\DeclareMathOperator{\Sym}{Sym}

\DeclareMathOperator{\Alb}{Alb}
\newcommand{\verti}{\mathrm{vert}}
\newcommand{\hor}{\mathrm{hor}}
\newcommand{\univ}{\mathrm{univ}}
\newcommand{\reg}{\mathrm{reg}}
\newcommand{\sing}{\mathrm{sing}}
\newcommand{\qt}{\mathrm{qt}}
\newcommand{\can}{\mathrm{can}}
\newcommand{\loc}{\mathrm{loc}}

\newcommand{\orb}{\mathrm{orb}}
\DeclareMathOperator{\tor}{Tor}
\newcommand{\Tor}{\mathrm{tor}}

\newcommand{\Sha}{\operatorname{Sha}}
\newcommand{\sha}{\operatorname{sha}}
\newcommand{\Shaf}{\mathrm{Sha}}
\newcommand{\shaf}{\mathrm{sha}}

% ============================================================
% Differential-geometric notation
% ============================================================
\newcommand{\ddbar}{\frac{\sqrt{-1}}{2\pi} \partial \bar{\partial}}
\newcommand{\deldel}{\sqrt{-1}\partial \overline{\partial}}
\newcommand{\dbar}{\overline{\partial}}

\newcommand{\pdrv}[2]{\frac{\partial #1}{\partial #2}}
\newcommand{\drv}[2]{\frac{d #1}{d#2}}
\newcommand{\ppdrv}[3]{\frac{\partial #1}{\partial #2 \partial #3}}

\newcommand{\polar}{\Omega}

% ============================================================
% Sheaves, ideals, automorphisms, kernels, and images
% ============================================================
\newcommand{\sO}{\mathcal{O}}
\newcommand{\mO}{\mathcal{O}}
\newcommand{\cV}{\mathcal{V}}
\newcommand{\I}[1]{\mathcal{I}(#1)}
\newcommand{\Aut}[1]{\Autop(#1)}
\newcommand{\Ker}[1]{\Kerop(#1)}
\newcommand{\Image}[1]{\Imageop(#1)}

% ============================================================
% Number systems and standard spaces
% ============================================================
\newcommand{\R}{\mathbb{R}}
\newcommand{\Z}{\mathbb{Z}}
\newcommand{\N}{\mathbb{N}}
\newcommand{\C}{\mathbb{C}}
\newcommand{\Q}{\mathbb{Q}}
\newcommand{\D}{\mathbb{D}}
\newcommand{\mP}{\mathbb{P}}
\newcommand{\B}{\mathds{B}}
\newcommand{\T}{\mathbb{T}}
\newcommand{\G}{\mathbb{G}}

% ============================================================
% Chern character, Todd class, Chow groups, and volumes
% ============================================================
\DeclareMathOperator{\ch}{ch}
\DeclareMathOperator{\td}{td}
\DeclareMathOperator{\Td}{Td}
\DeclareMathOperator{\CH}{CH}
\DeclareMathOperator{\vol}{vol}
\DeclareMathOperator{\Amp}{Amp}
\DeclareMathOperator{\Cone}{Cone}
\newcommand{\dR}{\mathrm{dR}}

% ============================================================
% Special symbols and formatting shortcuts
% ============================================================
%\newcommand{\tl}{\hspace{-0.8ex}<\hspace{-0.8ex}}
%\newcommand{\tr}{\hspace{-0.8ex}>}

\newcommand{\xb}[1]{\textcolor{blue}{#1}}
\newcommand{\xr}[1]{\textcolor{red}{#1}}
\newcommand{\xm}[1]{\textcolor{magenta}{#1}}

% This command places a short explanation below an equality or inequality sign.
% Example: A \underalign{\text{by Lemma 1.1}}{=} B.
\newcommand{\underalign}[2]{\quad \underset{\mathclap{\strut #1}}{#2}\quad}



\usepackage{booktabs,longtable}
\newcommand{\BC}{\mathrm{BC}}
\newcommand{\bas}{\mathrm{B}}
\newcommand{\virt}{\mathrm{virt}}
\newcommand{\Hol}{\operatorname{Hol}}
\newcommand{\SU}{\operatorname{SU}}
\newcommand{\Sp}{\operatorname{Sp}}
\newcommand{\Pic}{\operatorname{Pic}}
\newcommand{\ddc}{dd^{c}}
\newcommand{\wt}{\omega_{T}}
\newcommand{\rhoT}{\rho_{T}}
\newcommand{\rhoC}{\rho_{C}}
\newcommand{\AIstatus}[1]{\textnormal{\textsf{[#1]}}}
\newcommand{\Qfac}{\mathcal{Q}}
\emergencystretch=2em
\raggedbottom
\hypersetup{urlcolor=blue}

% BEGIN integration support (no manuscript prose)
% Author capitalization is the sole change to title typography.
\usepackage{etoolbox}
% Write bookmarks while each language namespace is active.
\usepackage{bookmark}
\makeatletter
\patchcmd{\@setauthors}{\MakeUppercase{\authors}}{\authors}{}{\errmessage{Author patch failed}}
\apptocmd{\maketitle}{\markboth{chatGPT 6 Astra}{\shorttitle}}{}{\errmessage{Author header patch failed}}
\let\MergeSetAbstract\@setabstract
\let\MergeSectionNumber\thesection
\let\MergeSectionName\sectionname
\let\MergeHSection\theHsection
\let\MergeChapterApp\Hy@chapapp
\newcommand{\MergeLanguage}{en}
% Namespaces affect internal lookup and PDF destinations, never source arguments.
\renewcommand*{\HyperDestNameFilter}[1]{\MergeLanguage.#1}
\renewcommand*{\pdfbookmark}[3][0]{%
  \bookmark[level=#1,dest={\HyperDestNameFilter{#3.#1}}]{#2}%
  \hyper@anchorstart{#3.#1}\hyper@anchorend}
\AtBeginDocument{%
  \let\MergeOriginalLabel\label
  \NewCommandCopy\MergeOriginalRef\ref
  \let\MergeOriginalCite\@citex
  \let\MergeOriginalBibitem\bibitem
  \RenewDocumentCommand\label{m}{\MergeOriginalLabel{\MergeLanguage:#1}}%
  \let\ltx@label\label % amsmath display labels use this saved entry point.
  \RenewDocumentCommand\ref{s m}{\IfBooleanTF{#1}{\MergeOriginalRef*{\MergeLanguage:#2}}{\MergeOriginalRef{\MergeLanguage:#2}}}%
  \def\@citex[#1]#2{%
    \def\MergeKeys{}%
    \@for\MergeKey:=#2\do{%
      \ifx\MergeKeys\@empty
        \edef\MergeKeys{\MergeLanguage:\MergeKey}%
      \else\edef\MergeKeys{\MergeKeys,\MergeLanguage:\MergeKey}\fi}%
    \MergeOriginalCite[#1]{\MergeKeys}}%
  \RenewDocumentCommand\bibitem{o m}{%
    \IfNoValueTF{#1}{\MergeOriginalBibitem{\MergeLanguage:#2}}{\MergeOriginalBibitem[#1]{\MergeLanguage:#2}}}%
  % Snapshot original counter values; page numbers intentionally remain continuous.
  \def\MergeResetCounters{}%
  \begingroup
    \def\@elt#1{\ifstrequal{#1}{page}{}{%
      \xappto\MergeResetCounters{\noexpand\setcounter{#1}{\number\csname c@#1\endcsname}}}}%
    \cl@@ckpt
  \endgroup
}
% Two independent Contents files, with original heading and formatting.
\newswitch{entoc}
\newswitch{jatoc}
\let\MergeOriginalStarttoc\@starttoc
\def\@starttoc#1#2{\MergeOriginalStarttoc{\MergeLanguage#1}{#2}}
\let\MergeOriginalAddtocontents\addtocontents
\renewcommand\addtocontents[2]{\MergeOriginalAddtocontents{\MergeLanguage#1}{#2}}
\newcommand{\MergeStart}[1]{%
  \def\MergeLanguage{#1}%
  \MergeResetCounters
  \let\thesection\MergeSectionNumber
  \let\sectionname\MergeSectionName
  \let\theHsection\MergeHSection
  \global\let\Hy@chapapp\MergeChapterApp
  \global\let\@setabstract\MergeSetAbstract
  \global\let\authors\@empty
  \gdef\shortauthors{chatGPT 6 Astra}%
  \global\let\addresses\@empty
}
\makeatother

% END integration support
\begin{document}

\begingroup
\MergeStart{en}
\title[Basic Chern vanishing on Vaisman manifolds]{Basic Chern vanishing and Beauville--Bogomolov structure in Vaisman geometry}
\author{chatGPT 6 Astra}







\date{8 September 2026, working note, version 0.1}
\renewcommand{\subjclassname}{\textup{2020} Mathematics Subject Classification}
\subjclass[2020]{Primary 53C55; Secondary 32Q25, 32M25, 57R18}
\keywords{Vaisman manifold, basic Chern class, Bott--Chern cohomology,
transverse Calabi--Yau geometry, elliptic orbi-bundle, virtual first Betti number}
\hypersetup{pdftitle={Basic Chern vanishing and Beauville--Bogomolov structure in Vaisman geometry},pdfauthor={chatGPT 6 Astra},pdfsubject={Expository working note; AI-generated; not human-verified}}


\pdfbookmark[0]{English version}{language-en}

\begin{abstract}
For a compact non-K\"ahler Vaisman manifold, vanishing of the first real
basic Chern class is equivalent to vanishing of the first Bott--Chern class
of the complex manifold. We explain this equivalence without assuming a
leaf space, and apply Istrati's existing structure theorem. The resulting
Beauville--Bogomolov decomposition belongs to the transverse geometry and
to a finite orbifold cover of the leaf space. The total space remains a
definite elliptic orbi-bundle. We compute its first and virtual first Betti
numbers, describe its fundamental group and Albanese map, and give explicit
examples with small first Betti number. Further elementary consequences
concern bundle reconstruction, product obstructions, and a vanishing second
basic Chern class. This is an expository working note; no new decomposition
theorem or verified claim of originality is asserted.
\end{abstract}
\maketitle
\xr{This note was generated by ChatGPT 6. Iwai has NOT verified any of its mathematical content. It may therefore contain mathematical errors and should be read with caution. A Japanese version is provided below.}
\tableofcontents

\section{The precise answer and the status of the results}
\label{sec:intro}
Throughout, $M$ is connected and compact, $\dim_{\C}M=n\geq2$, and
$(J,g,\omega,\theta)$ is Vaisman with $\theta\ne0$. Put $m=n-1$.
The hypothesis in this note is exactly
\begin{equation}\label{eq:hyp}
c_1^{\bas}(Q^{1,0})=0\quad\text{in }H^2_{\bas}(M,\R).
\end{equation}
It imposes no condition on $c_2^{\bas}$. The complex structure $J$ is fixed.

\begin{thm}[Answer to the main question]\label{thm:main}
Under \eqref{eq:hyp}, the original canonical foliation is quasi-regular.
There is a connected finite holomorphic \emph{\'etale} cover $M'\to M$ and
a holomorphic principal elliptic orbi-bundle
\begin{equation}\label{eq:BB}
E\longrightarrow M'\xrightarrow{\pi}B,
\qquad B=A\times\prod_i Y_i\times\prod_j Z_j=A\times\Qfac.
\end{equation}
Here $A$ is an abelian variety of dimension $a$, and the other factors are
projective pure K\"ahler orbifolds, simply connected in the \emph{orbifold}
sense, with Ricci-flat metrics and irreducible holonomy respectively
$\SU(d_i)$ and $\Sp(e_j)$. Their complex dimensions are $d_i$ and $2e_j$.
In particular,
\begin{equation}\label{eq:dimension}
a+\sum_i d_i+\sum_j2e_j=m.
\end{equation}
The product decomposition is holomorphic and isometric on the base for
the induced Ricci-flat metric. The total space is generally twisted; indeed
no finite \'etale cover of $M$ is a product of two positive-dimensional
compact complex manifolds. Furthermore,
\begin{equation}\label{eq:mainBetti}
b_1(M')=b_1^{\virt}(M)=2a+1,
\qquad b_1(M)\leq2a+1\leq2n-1.
\end{equation}
After choosing origins, $\Alb(M')=A$ and its Albanese map is
$\pr_A\circ\pi$. Its fibers have dimension $n-a$, and are the canonical
leaves precisely when $\Qfac$ is a point.
\end{thm}

The equivalence of Chern conditions and the structure assertion are
\AIstatus{KNOWN}: see Istrati \cite[Definition 4.1, Lemma 4.4,
Corollary 5.2, Theorem 5.3]{Istrati}. We supply the bridge and the
applications in detail. The virtual Betti formula and bundle computations
are \AIstatus{DERIVED} here, which is not a claim of novelty.
The integral fundamental-group descriptions below refine the explanation
of \cite[Corollary 5.10]{Istrati}. The bound and maximal case already appear
in \cite[Corollary 5.11]{Istrati}.

\begin{rem}[Version and notation policy]\label{rem:versions}
All numbered references to Istrati mean arXiv:2304.02582v1 (5 April 2023).
The publisher confirms JEMS online first, 11 March 2025,
DOI 10.4171/JEMS/1633. Its subscription full text was not accessible in this
run; a correspondence with published numbering is therefore not asserted.
Gomes means arXiv:2604.04134v2 (21 May 2026), including its revised title.
His real dimension $2r+2$ corresponds to our $r=m=n-1$; his maximal value
$2r+1$ is our $2n-1$. His Theorems 1.1 and 6.6 give a classification with
left-invariant complex structure and identify the leaf fibration with the
Albanese map. His proof uses a deformation of $J$, with a separate return
argument in Theorem 4.3. We do not transfer a statement from a deformed
complex structure: our route uses Istrati directly on $(M,J)$.
\end{rem}

\section{Definitions and transverse Hodge theory}\label{sec:defs}
We use $\omega(X,Y)=g(JX,Y)$ and set $J\alpha=-\alpha\circ J$ for a real
one-form. Multiplying $g$ by the positive constant $|\theta|_g^2$ makes
$|\theta|=1$, without changing $J$, $\theta$, or its Levi-Civita
parallelness. With
\begin{equation}\label{eq:vaisman}
U=\theta^\sharp,\quad V=JU,\quad\eta=J\theta,
\quad\wt=-d\eta,\quad \omega=\wt+\theta\wedge\eta,
\end{equation}
the form $\wt$ is closed, semipositive, of complex rank $m$, and its kernel
is $T\cF=\langle U,V\rangle_{\R}$. This is the transverse K\"ahler
structure. These standard Vaisman identities are recalled in
\cite[Section 2.1]{Istrati}, \cite[Theorem 2.2]{IO}.

\begin{defn}\label{def:basic}
A differential form $\alpha$ is basic if $\iota_X\alpha=\cL_X\alpha=0$
for every vector field $X$ tangent to $\cF$. Equivalently it is horizontal
and invariant for $U,V$. The complex $(\Omega_{\bas}^{\bullet},d_{\bas})$
has cohomology $H_{\bas}^{\bullet}(M,\R)$. Its complexification decomposes
by transverse type and $d_{\bas}=\partial_{\bas}+\bar\partial_{\bas}$.
We set
\begin{align}\label{eq:cohomdef}
H_{\bas}^{p,q}&=\frac{\ker(\bar\partial_{\bas}:\Omega_{\bas}^{p,q}
\to\Omega_{\bas}^{p,q+1})}{\bar\partial_{\bas}\Omega_{\bas}^{p,q-1}},\\
H_{\BC,\bas}^{1,1}(M,\R)&=
\frac{\{\alpha\in\Omega_{\bas}^{1,1}:d\alpha=0,\ \bar\alpha=\alpha\}}
{\{i\partial_{\bas}\bar\partial_{\bas}f:f\in C^\infty_{\bas}(M,\R)\}}.
\notag
\end{align}
Replacing basic forms and functions by ordinary ones defines
$H_{\BC}^{1,1}(M,\R)$.
\end{defn}

The holomorphic normal bundle is $Q^{1,0}=(TM/T\cF)^{1,0}$ of rank $m$.
In a transverse holomorphic chart, its Hermitian metric $h_T$ has Chern
connection matrix $h_T^{-1}\partial_{\bas}h_T$ and curvature $F_T$.
The invariant local matrices glue as a basic connection, equivalently a
connection extending the leafwise Bott partial connection. Define
\begin{equation}\label{eq:chern}
c^{\bas}(Q^{1,0})=
\left[\det\left(I+\frac{i}{2\pi}F_T\right)\right]_{\bas},
\quad \rhoT=i\tr F_T=-i\partial_{\bas}\bar\partial_{\bas}\log\det h_T.
\end{equation}
Changes of basic Hermitian metric change the Chern forms by basic exact
forms. For $TM$, the same convention gives
$\rhoC=-i\partial\bar\partial\log\det g$ and
$c_1^{\BC}(M)=[\rhoC/(2\pi)]_{\BC}$.
Thus we use Chern classes of \emph{tangent} bundles; the canonical bundle
has their negative first class. Our template macro
$\ddbar$ includes $1/(2\pi)$, whereas $i\partial\bar\partial$ does not.

\begin{lem}[The basic potential lemma]\label{lem:ddbar}
The canonical foliation is a transversely oriented Riemannian K\"ahler
foliation with minimal leaves. The basic $\partial\bar\partial$ lemma
applies. In particular, a real, basic, $d_{\bas}$-exact $(1,1)$-form
$\alpha$ can be written $\alpha=i\partial_{\bas}\bar\partial_{\bas}f$
with $f$ real and basic.
\end{lem}
\begin{proof}
The fields $U,V$ commute and are real holomorphic Killing fields of unit
length. Since $U$ is parallel, $\nabla_UU=\nabla_VU=0$; commutation gives
$\nabla_UV=0$, and the constant length Killing identity gives
$\nabla_VV=-\tfrac12\operatorname{grad}|V|^2=0$. The leaves are totally
geodesic, so their mean curvature vanishes. The metric is bundle-like,
and $\wt^m$ orients the normal bundle. Basic Hodge theory and the K\"ahler
identities therefore apply: \cite[Sections 3.3--3.4, Proposition 3.5.1]{EK},
\cite[Section 2.2]{IO}. These hypotheses are satisfied before there is any
global quotient orbifold.

For clarity, the required degree-two consequence has the following
Hodge-theoretic proof. Write $\chi=\theta\wedge\eta$; then
$d\chi=\theta\wedge\wt$. Integrating a basic exact form against
$\wt^{m-1}\wedge\chi$ gives zero by Stokes: the extra term contains
a transverse form of degree $2m+1$, so vanishes. In particular,
the trace $\Lambda_T\alpha$ has mean zero. The basic scalar Laplacian,
equivalently $f\mapsto\Lambda_T(i\partial_{\bas}\bar\partial_{\bas}f)$
up to a fixed nonzero scalar, is self-adjoint and has kernel the constants.
Compact basic Hodge theory solves
\[
\Lambda_T(i\partial_{\bas}\bar\partial_{\bas}f)=\Lambda_T\alpha.
\]
The difference $\beta$ is closed and primitive. For $m\geq2$ its
transverse Hodge star is
$*_T\beta=-\beta\wedge\wt^{m-2}/(m-2)!$, so it is also basic
coclosed. Since it is basic exact, its squared norm is zero. For $m=1$,
a trace-free $(1,1)$-form is already zero. This proves the assertion.
The same basic Hodge theory gives $H_{\bas}^{2m}(M,\R)=\R$; this is
the homological orientability required in El Kacimi-Alaoui's formulation.
\end{proof}

\begin{defn}\label{def:orb}
The canonical foliation is \emph{quasi-regular} when its leaves are compact,
or equivalently its effective Lee group is a compact complex
one-dimensional group. It is \emph{regular} when this group acts freely.
For a quasi-regular foliation the quotient $\cB$ is a complex orbifold.
A holomorphic principal $E$-orbi-bundle is locally of the form
$(\widetilde W\times E)/\Gamma\to\widetilde W/\Gamma$, with holomorphic
transition functions acting by translations of $E$; the isotropy acts on
the fiber by translations. The diagonal action is free when the total
space is a manifold. An ordinary principal bundle has a smooth base and
no orbifold isotropy.

A pure orbifold has no divisorial isotropy; it can still have quotient
singularities in codimension at least two. We write
$H^k_{\orb}(\cB,\Lambda)=H^k(B\cB,\Lambda)$ and
$\pi_k^{\orb}(\cB)=\pi_k(B\cB)$ for its classifying space.
Real orbifold cohomology agrees with orbifold de Rham cohomology and the
real cohomology of the underlying space. Integral cohomology need not.
An orbifold covering is locally a change from a finite chart group to
a subgroup. It need not be an \'etale cover of the underlying singular
spaces. By a finite \'etale cover of manifolds we mean a finite unramified
holomorphic covering of the total spaces.
\end{defn}

The LCK rank is the rank of the period group of $\theta$ (equivalently
the smallest rational subspace containing $[\theta]$). Rank one is not
the definition of either regularity or quasi-regularity. No rank-one
assumption is used in our structure theorem.

\section{The bridge to Bott--Chern vanishing}\label{sec:bridge}
\begin{lem}[Ricci comparison; Istrati's Lemma 4.4]\label{lem:ricci}
For a normalized Vaisman metric, as ordinary forms on $M$,
\begin{equation}\label{eq:ricci}
\rhoC=\rhoT.
\end{equation}
\end{lem}
\begin{proof}
Take a transverse holomorphic submersion $p:W\to T$ and a local
holomorphic frame $\sigma_T$ of $K_T$. The form
$\theta^{1,0}=(\theta+i\eta)/2$ is generally \emph{not} holomorphic.
Nevertheless
\[
\sigma=\theta^{1,0}\wedge p^*\sigma_T
\]
is a nonvanishing holomorphic frame of $K_M|_W$:
$\bar\partial\theta^{1,0}=-i\wt/2$, whose wedge with the horizontal
top form is zero. Its squared norm is
$|\sigma|_g^2=\tfrac12p^*|\sigma_T|_{h_T}^2$.
For a holomorphic canonical frame the tangent Ricci form is
$i\partial\bar\partial\log|\sigma|^2$. The constant $1/2$ disappears
upon differentiation, proving \eqref{eq:ricci}. This also checks the sign
and the absence of an additional multiple of $\wt$.
\end{proof}

\begin{prop}[Comparison of the three conditions]\label{prop:bridge}
For every compact non-K\"ahler Vaisman manifold,
\begin{equation}\label{eq:equiv}
c_1^{\bas}(Q^{1,0})=0
\ \Longleftrightarrow\ c_1^{\BC}(M)=0
\ \Longrightarrow\ c_1(M)=0\text{ in }H^2(M,\R).
\end{equation}
The last implication cannot be reversed.
\end{prop}
\begin{proof}
There is a commuting square of natural maps
\[
\begin{array}{ccc}
H_{\BC,\bas}^{1,1}(M,\R)&\xrightarrow{j}&H_{\BC}^{1,1}(M,\R)\\
\big\downarrow&&\big\downarrow\\
H_{\bas}^{2}(M,\R)&\longrightarrow&H^{2}(M,\R)
\end{array}
\]
The left map is injective by Lemma~\ref{lem:ddbar}, and its image is the
real transverse $(1,1)$-part. If $c_1^{\bas}=0$, that lemma applied to
$\rhoT$ gives $\rhoT=i\partial_{\bas}\bar\partial_{\bas}f$.
A basic function is an ordinary global smooth function, and the restricted
operators agree with $\partial,\bar\partial$ on it.
Lemma~\ref{lem:ricci} therefore proves Bott--Chern vanishing.

The reverse direction also holds before quasi-regularity is known.
Let $K$ be the closure of the flows of $U,V$ inside the compact isometry
group. It consists of holomorphic isometries. If a basic real $(1,1)$-form
$\alpha=i\partial\bar\partial f$, average $f$ over $K$ with Haar
probability measure. Since $\alpha$ is invariant, the average satisfies
$i\partial\bar\partial\bar f=\alpha$ and is invariant under $U,V$,
hence basic. This proves injectivity of $j$, and applied to $\rhoT=\rhoC$
proves the converse in \eqref{eq:equiv}. The last map is simply forgetting
Bott--Chern information.
\end{proof}

Here is an explicit description of the information that is forgotten.
\begin{lem}\label{lem:lowdegree}
For any compact non-K\"ahler Vaisman manifold,
\begin{align}\label{eq:lowdegree}
H^1(M,\R)&=H^1_{\bas}(M,\R)\oplus\R[\theta],\\
\ker\bigl(H^2_{\bas}(M,\R)\to H^2(M,\R)\bigr)&=\R[\wt]_{\bas}.
\notag
\end{align}
In particular $[\wt]_{\bas}\ne0$, although $\wt=-d\eta$ on $M$.
\end{lem}
\begin{proof}
Averaging over $K$ does not change ordinary de Rham classes because its
elements are isotopic to the identity. An invariant one-form uniquely
has the shape $\beta+u\theta+v\eta$, with $\beta$ basic and $u,v$ basic
functions. Its derivative is
$d_{\bas}\beta+du\wedge\theta+dv\wedge\eta-v\wt$.
If this is basic, contraction with $U,V$ gives $du=dv=0$.
The class $[\wt]_{\bas}$ is nonzero: a hypothetical basic primitive
would, by Stokes as in Lemma~\ref{lem:ddbar}, force
$\int_M\wt^m\wedge\chi=0$, whereas the integrand is positive.
Thus closed invariant one-forms have $v=0$, proving the first formula.
The sum is direct because contraction of an invariant exact one-form
$df$ with $U$ is zero after averaging $f$. Applying the same argument to
a primitive of any ordinary exact basic two-form proves the second
formula. This is also consistent with \cite[Sections 3--5]{IO}.
\end{proof}

For the scalar Hopf manifold
$(\C^n\setminus\{0\})/\langle q\rangle$, $0<q<1$, the leaf space is
$\mP^{n-1}$ and $c_1^{\bas}=nH\ne0$, while $H^2(M,\R)=0$.
Thus it is a counterexample to the reverse of the last implication only;
it is \emph{not} an example satisfying our main hypothesis.
Equivalently, $c_1(M)=0$ permits a nonzero multiple of
$[\wt]_{\bas}$ in place of zero.

\section{Applying the known structure theorem}\label{sec:structure}
\begin{proof}[Proof of the structural part of Theorem~\ref{thm:main}]
Proposition~\ref{prop:bridge} supplies exactly the hypothesis in Istrati's
Definition 4.1. Compactness, the Vaisman condition, and nonzero Lee form
are already part of the setting. Istrati's Corollary 5.2 gives a
Chern--Ricci-flat Vaisman metric on the same complex manifold, in a
prescribed admissible Lee class, unique up to scale. By
Lemma~\ref{lem:ricci} its transverse metric is Ricci-flat.

Theorem 5.3 of that paper proves, without a preliminary quotient
assumption, that $\Autop_0(M)$ is a compact complex one-dimensional group
and equals the Lee group. Its universal complex Lie group is $\C$ and
its kernel is a cocompact rank-two lattice; hence it is an elliptic curve
$E=\C/\Lambda$. The locally free action gives the original foliation
its compact leaves. Its quotient is now available. This is the step
that proves quasi-regularity and avoids a circular use of a K\"ahler
orbifold quotient.

For orientation, the mechanism in that theorem is that the uniqueness
of the canonical metric makes every element of $\Autop_0(M)$ an
isometry: it preserves the Lee class and total volume. The compact
automorphism-group result then identifies it with the Lee group. The
rest of Istrati's theorem applies an orbifold splitting theorem and
proves purity after a finite cover using the definite Seifert bundle.
For pure orbifolds, the decomposition used is
\cite[Theorem 6.4 and Definition 6.5]{Campana}. Its hypotheses are
compactness, K\"ahler geometry and vanishing orbifold first Chern class.
Here the last follows by identifying basic forms with quotient
orbifold forms and using \eqref{eq:hyp}. We use Istrati's purity argument
as part of the known theorem, rather than postulating purity initially.

The quotient has an integral positive class coming from the definite
bundle; hence it is projective. Every factor of its finite product
cover is projective by restriction of an ample class. Thus the flat
complex torus is an abelian variety, not just an arbitrary complex
torus. The meaning of simple connectivity is
$\pi_1^{\orb}(Y_i)=\pi_1^{\orb}(Z_j)=0$. Holonomy is that of the
Ricci-flat metric on the smooth locus of the pure orbifold, as in
Campana's definition. We label the $\SU$ factors by $Y_i$, reversing
the letters used for the two types in Istrati's Theorem 5.3.
The overlap $\SU(2)=\Sp(1)$ may be assigned to the symplectic type once;
it is not counted twice.

Finally, form the orbifold fiber product with this finite covering of
the quotient. Locally it is induced from
$(\widetilde W\times E)/\Gamma'\to(\widetilde W\times E)/\Gamma$,
where $\Gamma'\subset\Gamma$. Since the diagonal action on the original
total space is free, this is a genuine unramified covering of manifolds.
The complex structures and Vaisman metrics pull back. Taking a connected
component if necessary yields $M'$. This proves \eqref{eq:BB} and
\eqref{eq:dimension}. Neither regularity nor smoothness of the factors
has been asserted. Istrati's Corollary 5.4 also gives that $K_M$ is
holomorphically torsion.
\end{proof}

The transverse Calabi--Yau input can equivalently be expressed as follows:
for a compact homologically orientable transversely K\"ahler foliation,
each real basic $(1,1)$ representative of $2\pi c_1^{\bas}$ is the
Ricci form of a unique transverse K\"ahler metric in a fixed basic
K\"ahler class \cite[Section 3.5.5]{EK}. In transverse dimension one the
statement reduces to the scalar elliptic equation. The fixed-$J$
Vaisman implementation is \cite[Theorem 4.1, Lemma 4.2]{OV} and
Istrati's Corollary 5.2. We do not assume that an arbitrary conformal
change to a Chern--Ricci-flat Hermitian metric is still Vaisman.

\section{Betti numbers, finite covers, and the Albanese map}\label{sec:betti}
\begin{prop}[Transgression calculation]\label{prop:betti}
For the cover in \eqref{eq:BB},
\begin{equation}\label{eq:betti}
b_1(M')=2a+1.
\end{equation}
\end{prop}
\begin{proof}
A connection on the principal $E$-orbi-bundle gives its characteristic
class $c\in H^2_{\orb}(B,\Lambda)$. The Vaisman connection has
curvature proportional to $\wt$ in a single nonzero direction of
$\Lambda\otimes\R$. One way to see this is to use $\theta,\eta$
on the two-dimensional vertical tangent space: $d\theta=0$ and
$d\eta=-\wt$. Choosing a period basis only applies an invertible real
linear change to these two curvature components. Hence their span in
$H^2_{\orb}(B,\R)$ has rank one, not zero. Positivity and compactness
give $[\wt]^m\ne0$ on the base.

The real Leray--Serre sequence of classifying spaces has its degree-one
exact part
\begin{equation}\label{eq:five}
0\longrightarrow H^1_{\orb}(B,\R)\longrightarrow H^1(M',\R)
\longrightarrow H^1(E,\R)\xrightarrow{d_2}H^2_{\orb}(B,\R).
\end{equation}
There is no local coefficient monodromy on $H^1(E)$ because transition
functions are translations. If $u,v$ are the dual period basis,
$d_2(u),d_2(v)$ are the two characteristic classes, up to the common
choice of sign for transgression. Thus $\rank d_2=1$ and
$\dim\ker d_2=1$. Orbifold simple connectivity gives
$H^1_{\orb}(\Qfac,\R)=0$, so $b_1^{\orb}(B)=2a$. Taking dimensions
in \eqref{eq:five} proves the formula.
\end{proof}

\begin{prop}[Virtual first Betti number]\label{prop:virtual}
Define $b_1^{\virt}(M)$ as the supremum of $b_1(N)$ over connected finite
holomorphic \'etale covers $N\to M$. Then the supremum is achieved and
equals $2a+1$. In particular,
\begin{equation}\label{eq:gap}
b_1^{\virt}(M)\in\{1,3,\ldots,2n-5\}\cup\{2n-1\},
\end{equation}
where the first set is empty for $n=2$. The value $2n-3$ is excluded
for the virtual invariant, but need not be excluded for $b_1(M)$.
\end{prop}
\begin{proof}
On any finite cover, the pulled-back transverse Ricci-flat metric has
the same restricted transverse holonomy. The basic Bochner identity,
valid here because mean curvature is zero, gives for a basic harmonic
one-form $\alpha$
\[
0=\|\nabla^T\alpha\|_{L^2}^2+
\int\Ric^T(\alpha^\sharp,\alpha^\sharp)\,dv
=\|\nabla^T\alpha\|_{L^2}^2.
\]
Thus basic harmonic one-forms are parallel. The nonflat irreducible
$\SU$ and $\Sp$ factors have no invariant covectors. The space of
parallel transverse covectors on any finite cover of $M'$ therefore
has dimension at most $2a$. Lemma~\ref{lem:lowdegree} bounds its ordinary
first Betti number by $2a+1$. For a finite cover of $M$, take a connected
component of its fiber product with $M'$; transfer makes pullback in
real cohomology injective. The same bound follows. It is attained by
$M'$ by Proposition~\ref{prop:betti}.

Every nonflat factor has complex dimension at least two. A one-dimensional
Ricci-flat K\"ahler orbifold factor would be flat and belong to $A$;
it cannot be a positive-dimensional simply connected factor of this
decomposition. Hence either $a=m$ or $a\leq m-2$, proving
\eqref{eq:gap}. The Bochner step is also the mechanism behind
\cite[Proposition 6.4]{Gomes}, with its foliation hypotheses supplied here.
\end{proof}

\begin{prop}[Albanese description]\label{prop:alb}
One has $\dim\Alb(M)=(b_1(M)-1)/2$ and
$\Alb(M')\cong A$, with Albanese map $\pr_A\pi$ up to translation.
The Albanese map of $M$ is a surjective submersion. In general it is
not the quotient by the canonical foliation.
\end{prop}
\begin{proof}
For a compact Vaisman manifold, holomorphic one-forms are closed and
basic, and are exactly the basic harmonic forms of type $(1,0)$;
the period quotient is a genuine complex torus
\cite[Sections 2--3, Theorem 3.3]{Gomes}, \cite[Section 3]{IO}.
Basic Hodge symmetry and Lemma~\ref{lem:lowdegree} give the dimension.
On $M'$, these forms are exactly pullbacks of holomorphic one-forms
of $A$: the simply connected factors have none, by Ricci-flat Bochner
theory. The map $\pi_1(M')\to\pi_1^{\orb}(B)\to\pi_1(A)$ is
surjective because the fiber and $\Qfac$ are connected. Consequently
the period lattices, not only their rational spans, agree. The universal
property of the Albanese map now identifies it with $\pr_A\pi$.
Its fibers are the connected elliptic orbi-bundles over $\Qfac$.

For $M$ itself, using its Chern--Ricci-flat Vaisman metric, a basis of
holomorphic one-forms consists of transverse parallel forms. Their
pointwise independence follows from parallelness and independence at
one point. Thus the Albanese map has full target rank everywhere. Its
image is open by submersivity and closed by compactness, hence is the
whole connected Albanese torus.
\end{proof}

If $\widehat M\to M$ is a Galois cover with group $F$, then
$H^1(M,\R)=H^1(\widehat M,\R)^F$. This is the precise second source
of a small $b_1$: the finite group can remove transverse parallel forms.
It must be distinguished from the reduction in $a$ caused by $\Qfac$.
The statement $b_1(M)=2n-1$ if and only if $M$ is a Kodaira manifold
is \cite[Corollary 5.11]{Istrati}; the regular Albanese fibration and
left-invariant complex description are \cite[Theorems 1.1, 6.6]{Gomes}.
Without \eqref{eq:hyp} there is no dimension-only bound $2n-1$:
the circle-bundle construction of Section~\ref{sec:examples} over a
genus-$g$ curve gives Vaisman surfaces with $b_1=2g+1$ for all $g\geq2$.

\section{Characteristic classes and the limits of product decomposition}
\label{sec:twist}
For $E=\C/\Lambda$, principal holomorphic elliptic orbi-bundles are
classified by $H^1_{\orb}(B,\cO_B(E))$. The sequence
\begin{equation}\label{eq:exp}
0\longrightarrow\Lambda\longrightarrow\cO_B\longrightarrow
\cO_B(E)\longrightarrow0
\end{equation}
defines $c(M')\in H^2_{\orb}(B,\Lambda)$. Thus this characteristic
class is only part of the data: the fiber over a fixed class is a torsor
under the image of $H^1(B,\cO_B)$, namely the quotient by the image
of $H^1_{\orb}(B,\Lambda)$. One must retain the lattice of $E$ and
the holomorphic torsor class.

\begin{prop}[The line-bundle model and reconstruction]\label{prop:reconstruct}
For the product cover $B$ in \eqref{eq:BB}, choose an orientation of a
primitive lattice element $e$ and a complementary element $f$. Then
\begin{equation}\label{eq:char}
c(M')=\kappa\otimes e,
\qquad \kappa\in H^2_{\orb}(B,\Z),\quad \kappa\text{ a K\"ahler class}.
\end{equation}
There is an ample holomorphic line orbi-bundle $L$ with $c_1(L)=\kappa$
and $q\in\C^*$, $|q|<1$, such that
\begin{equation}\label{eq:line}
M'\cong L^{\times}/\langle q\rangle.
\end{equation}
One can write
$L\cong\pr_A^*L_A\otimes\bigotimes_i\pr_{Y_i}^*L_i
\otimes\bigotimes_j\pr_{Z_j}^*L_j'$.
The elliptic torsor is the Baer sum of the pullbacks of the corresponding
factor torsors, using the same $E$ and $q$.
\end{prop}
\begin{proof}
The real curvature calculation in Proposition~\ref{prop:betti} has
one-dimensional image. Evaluating its integral class on integral
two-cycles shows that this direction contains a nonzero lattice vector;
its line is therefore rational. Taking a primitive vector and its sign
gives a positive class $\kappa$ modulo torsion. On the product cover,
$H_1^{\orb}(B,\Z)=\Z^{2a}$, so the universal coefficient theorem
makes $H^2_{\orb}(B,\Z)$ torsion-free. Consequently no torsion term
remains and \eqref{eq:char} holds integrally.

The orbifold exponential sequence and the K\"ahler $(1,1)$ theorem
give a line orbi-bundle $L_0$ with this first Chern class. Identify
$\Lambda=\Z+\tau\Z$, $\operatorname{Im}\tau>0$, sending $e$ to $1$.
Then $E=\C^*/\langle q\rangle$, $q=e^{2\pi i\tau}$, and the torsor
$L_0^{\times}/\langle q\rangle$ has the prescribed characteristic
class. By \eqref{eq:exp} its difference from $M'$ lies in the image
of $H^1(B,\cO_B)$. The ordinary exponential sequence for line
orbi-bundles maps that same $H^1(B,\cO_B)$ to $\Pic^0_{\orb}(B)$.
Tensoring $L_0$ by its image absorbs the difference. This proves
\eqref{eq:line} without assuming that a topological class alone
determines a holomorphic bundle. Positivity is unaffected by a
$\Pic^0$ twist. This is the setting of \cite[Theorem 2.4]{Istrati}.

For the factorization, integral K\"unneth in degree two has no cross
terms involving a simply connected factor. Thus $c_1(L)$ is the sum
of the pullbacks of its restrictions to the factors. The quotient of
$L$ by the tensor product of those restrictions has zero first class.
Also $H^1(B,\cO_B)=H^1(A,\cO_A)$, by K\"ahler Hodge theory and
$H^1_{\orb}(\Qfac,\C)=0$. Its remaining $\Pic^0$ term therefore
comes from $A$ and can be absorbed into $L_A$.

Explicitly, write $P_k\to B$ for the pulled-back factor torsors.
On their fiber product quotient by
\[
\ker(E^r\xrightarrow{\mathrm{sum}}E)
\]
acting by translations. Locally the fiber is $E^r/\ker(\mathrm{sum})=E$.
The resulting transition functions are the sums of the factor transition
functions, which in multiplicative $\C^*$ coordinates are the products
of the line-bundle transition functions. They are precisely those of
\eqref{eq:line}. This proves the reconstruction assertion.
\end{proof}

Before this cover a definite bundle can have integral torsion in addition
to its positive rank-one real class. Here torsion disappears because
the specific product cover has torsion-free $H_1$, not because every
torsion class on every orbifold can always be removed by an arbitrary
finite cover. The positive class never disappears under a finite base
cover: the real transfer identity $p_*p^*=(\deg p)\id$ makes pullback
injective. Fiber isogenies likewise do not annihilate it over $\R$.

\begin{prop}[A stronger obstruction to holomorphic products]\label{prop:noproduct}
A compact non-K\"ahler Vaisman manifold cannot be biholomorphic to
$X_1\times X_2$ with both $X_i$ compact and positive-dimensional.
The assertion persists under finite \'etale covers.
\end{prop}
\begin{proof}
Restrict the exact semipositive form $\wt$ to either compact slice
$X_i$, of dimension $r_i$. Stokes gives
$\int_{X_i}\wt^{r_i}=0$. Semipositivity then implies that the
restriction has deficient rank at every point. For a semipositive
Hermitian form, a vector of zero norm is in its full kernel. Hence
$T X_i\cap\ker\wt\ne0$ for both factors at every point of the
product. Since $\ker\wt$ is a complex line, it would have to lie in
both complementary tangent factors, a contradiction. A finite \'etale
cover is again Vaisman, so the argument repeats. This recovers the
product obstruction mentioned in \cite[Introduction]{Istrati}.
\end{proof}

In the line model, choose any Hermitian norm. Writing a nonzero vector
as its unit vector times its positive radius identifies $L^{\times}$
smoothly with $S(L)\times\R$. If $q$ has argument $\varphi$, the
quotient is a mapping torus of circle rotation by $\varphi$ on $S(L)$.
That rotation is isotopic to the identity, so
\begin{equation}\label{eq:smooth}
M'\simeq_{\mathrm{diff}} S(L)\times S^1.
\end{equation}
The unit-circle total space is smooth in our situation because the
isotropy acts freely on the nonzero line. For a suitable norm this
is the Sasaki construction. Formula \eqref{eq:smooth} is a real smooth
product and does not assert any holomorphic product. In particular
$E\times B$ would have first Betti number $2a+2$, different from
\eqref{eq:betti}. \emph{The BB analogy is a product theorem for the
base and transverse metric, and a bundle/extension theorem for the
total space.}

\section{The fundamental group}\label{sec:pi}
The homotopy exact sequence, interpreted on classifying spaces, is
\begin{equation}\label{eq:pi}
\pi_2^{\orb}(B)\xrightarrow{\langle c,\cdot\rangle}\Lambda
\longrightarrow\pi_1(M')\longrightarrow\Z^{2a}\longrightarrow1.
\end{equation}
The first arrow must not be omitted. In particular, an elliptic fiber
does not necessarily inject its fundamental group into the total space.

\begin{prop}\label{prop:pi}
If $\Qfac$ is a point, $\pi_1(M')$ is the product of $\Z$ with a
lattice in the real $(2a+1)$-dimensional Heisenberg group. It is not
virtually abelian. If $\dim\Qfac>0$, $\pi_1(M)$ is virtually
$\Z^{2a+1}$. These are the two cases in
\cite[Corollary 5.10]{Istrati}.
\end{prop}
\begin{proof}
When $B=A$, $\pi_2(B)=0$. A polarization basis puts $\kappa$ in
elementary-divisor form $(d_1,\ldots,d_a)$, all $d_i>0$. With a
choice of orientation of the central circle, a presentation is
\begin{equation}\label{eq:heis}
\langle x_i,y_i,z,t\mid z,t\text{ central},\
[x_i,y_j]=z^{d_i\delta_{ij}},\ [x_i,x_j]=[y_i,y_j]=1\rangle.
\end{equation}
The radial circle gives $t$. This presentation follows by lifting the
fundamental parallelograms of $A$: their circle winding numbers are
the evaluations of $c_1(L)$. The alternating commutator pairing is
nondegenerate over $\Q$, so no subgroup of finite index is abelian.

When $\Qfac$ has positive dimension, orbifold Hurewicz gives
$\pi_2^{\orb}(\Qfac)\cong H_2^{\orb}(\Qfac,\Z)$.
The restriction of the K\"ahler class $\kappa$ is nonzero, so the
image of the first arrow in \eqref{eq:pi} is $d\Z e$ for some $d>0$.
The kernel of the last arrow is therefore $\Z/d\oplus\Z$.
In the line model the $\Z$ factor splits off smoothly by
\eqref{eq:smooth}. The circle-bundle group is a central extension of
$\Z^{2a}$ by $\Z/d$, because the structure group consists of circle
translations. Its commutators are finite. The inverse image of
$d\Z^{2a}$ is abelian: commutators of lifts of multiples by $d$
are $d^2$-powers of commutators. It is finitely generated abelian,
so contains a torsion-free finite-index subgroup of rank $2a$.
Restoring the radial $\Z$ gives $\Z^{2a+1}$ up to finite index.
Passing from $M'$ to $M$ does not change this virtual conclusion.
\end{proof}

\section{Explicit examples and boundary cases}\label{sec:examples}
\begin{eg}[The construction used in the examples]\label{eg:construction}
Let $B$ be smooth projective, $H$ ample, $\kappa=c_1(H)$, and put
$L=H^{-1}$. Take a Hermitian metric $h$ on $L$ with
$i\partial\bar\partial\log h=2\pi\kappa_h>0$ and let
$F(v)=|v|_h^2$ on $L^{\times}$. Then
\begin{equation}\label{eq:cone}
\Omega_K=i\partial\bar\partial F>0,\qquad
\Omega_V=F^{-1}\Omega_K,\qquad\theta=-d\log F.
\end{equation}
Indeed, the horizontal term is $2\pi F\kappa_h$ and the vertical
term is $iF\partial\log F\wedge\bar\partial\log F$, both positive
in their respective directions. Thus $d\Omega_V=\theta\wedge\Omega_V$.
In the logarithmic radial coordinate its metric is a constant multiple
of a product $dt^2+g_{S(L)}$, so $\theta$ is parallel. Multiplication
by $q$, $0<q<1$, preserves $\Omega_V$ and $\theta$; the quotient
$M(H,q)=L^{\times}/\langle q\rangle$ is compact Vaisman.
The canonical leaves are its $E=\C^*/\langle q\rangle$ fibers,
and the quotient is $B$. In the counterclockwise fiber orientation,
\begin{equation}\label{eq:exchar}
c(M(H,q))=(-\kappa,0),\qquad
c_j^{\bas}(M(H,q))=c_j(TB).
\end{equation}
Replacing the first period generator by its negative gives the positive
convention of \eqref{eq:char}. The inversion $v\mapsto v^{-1}$
identifies the punctured bundles of $H$ and $H^{-1}$ with reversed
$\C^*$ action, so this sign choice agrees with the ample-line model.
The same local construction applies to a positive line orbi-bundle
whose nonzero total space is smooth.
\end{eg}

\begin{eg}[Maximal $b_1$: Kodaira manifolds]\label{eg:torus}
Let $B=A$ be a polarized abelian variety of dimension $m$, and use
Example~\ref{eg:construction}. The resulting manifold has dimension
$n=m+1$, is regular, has $c_j^{\bas}=0$ for every $j>0$, and
$b_1=2m+1=2n-1$. Its fundamental group is \eqref{eq:heis}, with
the sign of $z$ changed if necessary. Its Albanese map is the elliptic
bundle projection. These are exactly the models in
\cite[Example 7.2, Corollary 5.11]{Istrati}, compatible with
\cite[Theorem 6.6]{Gomes}. The standard model metric has LCK rank one;
this is not an assertion about every Vaisman metric on this complex
manifold. For $m=1$ it is a primary Kodaira surface.
\end{eg}

\begin{eg}[A K3 base and $b_1=1$]\label{eg:k3}
Let $S\subset\mP^3$ be a smooth quartic K3 surface and
$H=\cO_S(1)$. Adjunction gives $K_S\cong\cO_S$, and $S$ is
simply connected by the Lefschetz hyperplane theorem. Its hyperplane
class $h$ is primitive: $h^2=4$, and a possible divisibility $h=2u$
would give $u^2=1$, contradicting the even intersection form of a K3
surface. No larger divisibility is numerically possible. The intersection
pairing is unimodular by Poincar\'e duality, so evaluation of $h$ on
$H_2(S,\Z)$ is onto $\Z$.
For the K3 cohomology and lattice facts used here, see
\cite[Chapters 1 and 14]{Huybrechts}.

For $M=M(H,q)$, $n=3$, the canonical foliation is regular with base
$S$, and
\begin{equation}\label{eq:k3numbers}
c_1^{\bas}=0,\quad \int_S c_2^{\bas}=24,\quad
b_1(M)=b_1^{\virt}(M)=1,\quad \pi_1(M)=\Z.
\end{equation}
The $c_2$ calculation follows directly from
$c(TS)=(1+h)^4/(1+4h)$: its degree-two term is $6h^2$, whose
integral is $24$. The circle-bundle sequence
$\pi_2(S)\to\Z\to\pi_1(S(L))\to0$, together with primitivity,
gives $\pi_1(S(L))=0$. The radial circle supplies the displayed
fundamental group. The characteristic pair is $(-h,0)$.
The Albanese is a point, although the leaf space is a K3 surface.
Smoothly $M=S(L)\times S^1$; holomorphically it admits no nontrivial
compact product decomposition.
\end{eg}

\begin{eg}[Small $b_1$ caused by a finite group]\label{eg:secondary}
Take an elliptic curve $C$, the symmetric negative line bundle
$L=\cO_C(-[0])$, and the involution $\iota:x\mapsto-x$.
It has a line-bundle lift $\beta$ with $\beta^2=1$ (rescale any lift
by a constant square root). Choose a $\beta$-invariant negative
Hermitian metric and $0<q<1$. Set
\begin{equation}\label{eq:secondary}
P=L^{\times}/\langle q^2\rangle,
\qquad X=L^{\times}/\langle q\beta\rangle.
\end{equation}
The induced involution on $P$ is free. A fixed point would require
$q\beta(v)=q^{2k}v$ for some $k\in\Z$; taking norms gives
$q=q^{2k}$, which is impossible. The invariant metric descends, so
$P\to X$ is an \'etale double cover of Vaisman surfaces. On
$H^1(P,\R)=H^1(C,\R)\oplus\R[\theta]$, the involution acts as
$-1$ and $+1$ respectively. Therefore
\begin{equation}\label{eq:secondarynumbers}
b_1(X)=1,\qquad b_1(P)=b_1^{\virt}(X)=3.
\end{equation}
The leaf space of $X$ is the flat orbifold
$C/\{\pm1\}=\mP^1(2,2,2,2)$, so $c_1^{\bas}(X)=0$.
It is quasi-regular and not regular. Its characteristic class pulls
back to $(-[0],0)$ on $C$, and has nonzero rank-one real part;
the local order-two stabilizers act on the fiber by translation by
$q\beta|_{L_x}$ in $\C^*/\langle q^2\rangle$. This specifies the
isotropy data in addition to the pulled-back characteristic class.

For a compatible standard square-torus model, the group is
\[
\pi_1(X)=\langle x,y,z,s\mid z\text{ central},\ [x,y]=z,
\ sxs^{-1}=x^{-1},\ sys^{-1}=y^{-1},\ szs^{-1}=z\rangle.
\]
Equivalently it is $H_{\Z}\rtimes\Z$, with generator $s$ acting by
horizontal inversion. The subgroup generated by $x,y,z,s^2$ is
$H_{\Z}\times\Z$ of index two. One sees this model directly on the
real Heisenberg group with multiplication
$(x,y,z)(x',y',z')=(x+x',y+y',z+z'+xy')$: horizontal inversion
$(x,y,z)\mapsto(-x,-y,z)$ preserves its integral lattice and its
standard Sasaki structure. Translation by half the radial period makes
the quotient action free. This realizes the holomorphic construction
for the induced degree-one polarization. These are secondary Kodaira
examples, as in \cite[Example 7.3]{Istrati}.
\end{eg}

\begin{eg}[An orbifold factor that cannot be made smooth]\label{eg:orb}
The weighted degree-five polynomial
\[
f=z_0^2z_1+z_1^5+z_2^5+z_3^5,
\qquad (w_0,w_1,w_2,w_3)=(2,1,1,1)
\]
defines a smooth punctured hypersurface $W=\{f=0\}\setminus\{0\}$.
Indeed its four partial derivatives have common zero only at the
origin. Quotient by the weighted $\C^*$ action gives a projective
K3 orbifold $B\subset\mP(2,1,1,1)$, with an $A_1$ orbifold point
at $[1,0,0,0]$ and no divisorial isotropy. Adjunction gives trivial
orbifold canonical bundle since $5=2+1+1+1$. The quotient
$W/\langle q\rangle$ by weighted multiplication is a smooth compact
Vaisman threefold with $c_1^{\bas}=0$.

These are the explicit data of \cite[Example 7.4]{Istrati}. The link
of this isolated hypersurface singularity of complex dimension three
is simply connected by the hypersurface-link connectivity theorem
\cite[Chapter 5]{Milnor}. The sequence of its circle orbi-bundle
therefore gives $\pi_1^{\orb}(B)=0$. Hence $B$ has no nontrivial
connected orbifold covering at all, and in particular no finite smooth
orbifold cover. The total space has $\pi_1=\Z$ and $b_1=1$ from
the simply connected punctured cone and its infinite cyclic deck group.
Its characteristic pair is that of the tautological negative line
orbi-bundle $\cO_B(-1)$, namely $(-c_1^{\orb}(\cO_B(1)),0)$.
The local $\Z/2$ acts faithfully on the punctured tautological line.
This is an actual obstruction to replacing ``orbifold'' by ``manifold''
in the structure theorem, not merely an absence of such a proof.
\end{eg}

The three principal examples can be compared without conflating a
manifold and its cover:
\begin{center}\small
\begin{tabular}{lllll}\toprule
Example & $n$ & base & $b_1$ & $b_1^{\virt}$\\\midrule
Kodaira manifold & $m+1$ & $A^m$ & $2m+1$ & $2m+1$\\
K3 construction & $3$ & $S$ & $1$ & $1$\\
Secondary Kodaira & $2$ & $C/\{\pm1\}$ & $1$ & $3$\\\bottomrule
\end{tabular}
\end{center}

\section{Further consequences proved in this note}\label{sec:extra}
\begin{cor}[Attainable values]\label{cor:values}
For each fixed $n\geq2$, every value in \eqref{eq:gap} occurs as a
virtual first Betti number. The ordinary first Betti numbers occurring
under \eqref{eq:hyp} are exactly $1,3,\ldots,2n-1$.
\end{cor}
\begin{proof}
The virtual maximum is Example~\ref{eg:torus}. For $0\leq a\leq m-2$,
put $k=m-a\geq2$ and let $Q_k$ be a smooth degree-$k+2$ hypersurface
of $\mP^{k+1}$. Adjunction gives $K_{Q_k}\cong\cO$ and Lefschetz
gives simple connectivity. Its Ricci-flat metric has no flat factor,
since any such factor would give nonzero $H^1$. Apply
Example~\ref{eg:construction} to $A^a\times Q_k$ with an exterior
tensor-product polarization. Propositions~\ref{prop:betti} and
\ref{prop:virtual} give $b_1=b_1^{\virt}=2a+1$.

For the ordinary invariant choose a product $A=C^m$ of elliptic
curves and its product symmetric polarization. Let $\beta$ cover the
involution that is the identity on $\ell$ factors and $-1$ on the
remaining factors, $0\leq\ell<m$. The construction
\eqref{eq:secondary} works verbatim: radial contraction ensures
freeness even over fixed points. Its invariant transverse one-forms
have dimension $2\ell$, giving $b_1=2\ell+1$ and virtual invariant
$2m+1$. The case $\ell=m$ is the torus-base model itself.
Oddness follows in all cases from basic Hodge symmetry and
Lemma~\ref{lem:lowdegree}, and the upper bound was proved above.
\end{proof}

\begin{cor}[Adding the second basic Chern class]\label{cor:c2}
Under \eqref{eq:hyp}, the following are equivalent:
\begin{enumerate}
\item $c_2^{\bas}(Q^{1,0})=0$ in real basic cohomology;
\item there is a transverse flat K\"ahler metric for the fixed canonical
foliation and transverse complex structure;
\item $b_1^{\virt}(M)=2n-1$;
\item $M$ has a finite \'etale cover which is a Kodaira manifold.
\end{enumerate}
For $n=2$, the first condition is automatic and all four hold.
\end{cor}
\begin{proof}
Work on the quotient with its Ricci-flat metric, or on its finite
product cover. For $m\geq2$ its unitary Chern curvature $F$ has
$\tr F=0$ and $\Lambda F=0$. With the norm induced by $-\tr$ on
skew-Hermitian matrices, the primitive-form Hodge identity gives
\begin{equation}\label{eq:c2}
\int_B c_2(TB)\wedge\frac{\omega_B^{m-2}}{(m-2)!}
=\frac1{8\pi^2}\int_B |F|^2\frac{\omega_B^m}{m!}.
\end{equation}
Indeed $c_2=(8\pi^2)^{-1}\tr(F\wedge F)$ when $\tr F=0$,
and $*F=-F\wedge\omega_B^{m-2}/(m-2)!$. Both formulas apply in
orbifold charts with the usual division by stabilizer order in the
integral. Vanishing of the basic class therefore forces $F=0$.
In transverse dimension one, Ricci-flatness itself is flatness.
A compact flat K\"ahler orbifold has a finite torus orbifold cover
(the flat case of the splitting theorem used above); pulling it back
gives a Kodaira cover. Equivalently the nonflat factors in
\eqref{eq:BB} are absent. Proposition~\ref{prop:virtual} then gives
the third assertion. Conversely a Kodaira cover has zero basic Chern
classes, and real transfer for the lifted foliation is injective, so
the second basic class downstairs vanishes. The same decomposition
also proves directly that (3) implies (4).
\end{proof}

These consequences are proved here from existing structure and Hodge
theorems. The searches made for this note did not establish that they
are new in the literature; no originality claim is made.

\appendix
\section{Dependencies, checks, and what is not claimed}\label{sec:audit}
The logical route is: minimal transverse K\"ahler foliation;
basic potential lemma and Ricci comparison; equivalence
\eqref{eq:equiv}; Istrati's fixed-$J$ existence and quasi-regularity;
its finite orbifold BB decomposition; characteristic-class calculation;
Betti, Albanese, and fundamental-group consequences.

Compactness is used for Hodge theory, averaging, positivity integrals,
the automorphism group argument, and the Albanese image. The Vaisman
condition supplies minimal leaves, holomorphic Killing fields and
\eqref{eq:vaisman}. The first basic Chern condition enters precisely
at \eqref{eq:equiv} and transverse Ricci-flatness. Quasi-regularity is
only used after the known theorem proves it. Finite covers remove
transverse finite holonomy and the specified integral torsion; they do
not erase positive curvature classes. Orbifold structures are retained
in all homotopy, integral cohomology and covering arguments. Projectivity
comes from positivity and is used to choose the explicit algebraic
examples and factor polarizations.

The central question is answered by known results. No unresolved step
is deliberately used as a proved theorem here. We have not classified
all finite group actions on the factor torsors, or all isomorphisms
between their unmarked total spaces. These are further classification
problems, not missing hypotheses in Theorem~\ref{thm:main}.
No assertion that every orbifold is finitely smoothable is used.
No claim about arbitrary LCK manifolds follows from our proof.

The audit distinguishes \AIstatus{KNOWN}, \AIstatus{DERIVED},
\AIstatus{CANDIDATE}, \AIstatus{REFUTED}, and \AIstatus{GAP}.
Here the false implications ``de Rham vanishing implies basic
vanishing'', ``small $b_1$ forces a non-torus factor'', and ``a finite
cover makes every base smooth'' are \AIstatus{REFUTED} by explicit
examples. No \AIstatus{CANDIDATE} is presented as a theorem.
The inaccessible published Istrati numbering and the lack of an
exhaustive novelty search are bibliographic limitations, recorded in
\texttt{research\_audit\_ja.md}. Theorem numbers of the three requested
arXiv papers were checked in their actual bodies. The precise numbering
of the general hypersurface-link connectivity theorem in Milnor's
book was not checked; only its standard chapter-level attribution is
used. All references to deep inputs state the hypotheses needed here.
The separate README records compilation and visual verification.

\clearpage
% Bibliography recovered solely from the corresponding supplied PDF.
\begin{thebibliography}{EKA90}
\bibitem[Cam04]{Campana}
Frédéric Campana. Orbifoldes à première classe de Chern nulle. In \emph{The Fano Conference}, pages 339--351. Università di Torino, 2004. Theorem 6.4 and Definition 6.5 checked in arXiv:math/0402243v1.

\bibitem[EKA90]{EK}
Aziz El Kacimi-Alaoui. Opérateurs transversalement elliptiques sur un feuilletage riemannien et applications. \emph{Compositio Mathematica}, 73(1):57--106, 1990.

\bibitem[Gom26]{Gomes}
Lucas H. S. Gomes. On the classification of Vaisman manifolds with vanishing first basic Chern class and large first Betti number, 2026. arXiv:2604.04134v2, 21 May 2026.

\bibitem[Huy16]{Huybrechts}
Daniel Huybrechts. \emph{Lectures on K3 Surfaces}. Cambridge University Press, 2016. Author's online manuscript, Chapters 1 and 14 consulted.

\bibitem[IO23]{IO}
Nicolina Istrati and Alexandra Otiman. Bott--Chern cohomology of compact Vaisman manifolds. \emph{Transactions of the American Mathematical Society}, 376(6):3919--3936, 2023. Numbered citations refer to the final preprint arXiv:2206.07312v2, 26 September 2022.

\bibitem[Ist25]{Istrati}
Nicolina Istrati. Vaisman manifolds with vanishing first Chern class. \emph{Journal of the European Mathematical Society}, 2025. Published online first, 11 March 2025. Numbered citations in this note refer to arXiv:2304.02582v1, 5 April 2023.

\bibitem[Mil68]{Milnor}
John Milnor. \emph{Singular Points of Complex Hypersurfaces}, volume 61 of \emph{Annals of Mathematics Studies}. Princeton University Press, 1968. Chapter 5; precise theorem number for link connectivity not verified in this run.

\bibitem[OV24]{OV}
Liviu Ornea and Misha Verbitsky. A Calabi--Yau theorem for Vaisman manifolds. \emph{Communications in Analysis and Geometry}, 32(10):2889--2900, 2024. Numbered citations refer to arXiv:2206.08808v2, 24 August 2022.

\end{thebibliography}



% Language boundary: drain English floats before changing state.
\newpage
\clearpage
\endgroup

\begingroup
\MergeStart{ja}
\title[Vaisman多様体の第一basic Chern類と構造]{Vaisman幾何における第一basic Chern類の消滅とBeauville--Bogomolov型構造}
\author{chatGPT 6 Astra}







\date{8 September 2026, working note, version 0.1}
\renewcommand{\subjclassname}{\textup{2020} Mathematics Subject Classification}
\subjclass[2020]{Primary 53C55; Secondary 32Q25, 32M25, 57R18}
\keywords{Vaisman manifold, basic Chern class, Bott--Chern cohomology,
transverse Calabi--Yau geometry, elliptic orbi-bundle, virtual first Betti number}
\hypersetup{pdftitle={Vaisman幾何における第一basic Chern類の消滅とBeauville--Bogomolov型構造},pdfauthor={chatGPT 6 Astra},pdfsubject={Expository working note; AI-generated; not human-verified}}


\pdfbookmark[0]{日本語版}{language-ja}

\begin{abstract}
コンパクト非K\"ahler Vaisman多様体について，第一実basic Chern類の消滅は，
複素多様体自身の第一Bott--Chern類の消滅と同値である．本稿では，葉空間の存在を
仮定せずにこの同値を説明し，Istratiの既存の構造定理を適用する．得られる
Beauville--Bogomolov分解は横断幾何と葉空間の有限orbifold被覆に関するものであり，
全空間はdefiniteな楕円曲線orbi-bundleとして残る．第一Betti数とvirtual第一Betti数を
計算し，基本群とAlbanese写像を記述するとともに，第一Betti数が小さい具体例を与える．
さらに，束の再構成，積分解の障害，第二basic Chern類の消滅に関する初等的な帰結を示す．
本稿は解説的な作業ノートであり，新しい分解定理や確認済みの新規性を主張しない．
\end{abstract}
\maketitle
\xr{このノートはchatGPT 6が生成したものであり, 岩井は数学的な検証を全く行っていません. そのため数学的な間違いがあるかもしれません. そのため注意深く読んでください.}
\tableofcontents

\section{中心問題への正確な回答と結果の位置づけ}\label{sec:intro}
以下，$M$はコンパクト連結，$\dim_{\C}M=n\geq2$とし，
$(J,g,\omega,\theta)$は$\theta\ne0$を満たすVaisman構造とする．$m=n-1$とおく．
本稿の仮定は正確に
\begin{equation}\label{eq:hyp}
c_1^{\bas}(Q^{1,0})=0\quad\text{in }H^2_{\bas}(M,\R)
\end{equation}
である．$c_2^{\bas}$には何も仮定しない．複素構造$J$は固定する．

\begin{thm}[中心問題への回答]\label{thm:main}
\eqref{eq:hyp}のもとで，元のcanonical foliationはquasi-regularである．
連結な有限正則\emph{\'etale}被覆$M'\to M$と，正則principal elliptic orbi-bundle
\begin{equation}\label{eq:BB}
E\longrightarrow M'\xrightarrow{\pi}B,
\qquad B=A\times\prod_i Y_i\times\prod_j Z_j=A\times\Qfac
\end{equation}
が存在する．ここで$A$は次元$a$のabelian varietyであり，残りの因子は
\emph{orbifoldの意味で}単連結な射影的pure K\"ahler orbifoldである．
これらにはRicci-flat計量があり，既約holonomyはそれぞれ$\SU(d_i)$と$\Sp(e_j)$，
複素次元は$d_i$と$2e_j$である．特に
\begin{equation}\label{eq:dimension}
a+\sum_i d_i+\sum_j2e_j=m
\end{equation}
が成り立つ．底の分解は，誘導されるRicci-flat計量に関して，正則かつ等長的である．
全空間には束のねじれが残る．実際，$M$のどの有限\'etale被覆も，正の次元をもつ
二つのコンパクト複素多様体の積にはならない．さらに
\begin{equation}\label{eq:mainBetti}
b_1(M')=b_1^{\virt}(M)=2a+1,
\qquad b_1(M)\leq2a+1\leq2n-1
\end{equation}
が成り立つ．原点を選ぶと$\Alb(M')=A$で，Albanese写像は$\pr_A\circ\pi$である．
そのファイバーの次元は$n-a$であり，canonical foliationの葉と一致するのは
$\Qfac$が一点の場合に限る．
\end{thm}

Chern条件の同値と構造に関する主張は\AIstatus{KNOWN}である．
Istrati \cite[Definition 4.1, Lemma 4.4, Corollary 5.2, Theorem 5.3]{Istrati}
を参照されたい．本稿では条件の橋渡しと適用を詳しく説明する．virtual Betti数の式と
束の計算は本稿で\AIstatus{DERIVED}と分類するが，これは新規性の主張ではない．
後述の整係数での基本群の記述は\cite[Corollary 5.10]{Istrati}の説明を具体化する．
上界と最大値の場合は\cite[Corollary 5.11]{Istrati}ですでに得られている．

\begin{rem}[版と記号の規約]\label{rem:versions}
Istratiの定理番号はすべてarXiv:2304.02582v1，2023年4月5日のものを用いる．
出版社の情報から，JEMS online first，2025年3月11日公開，
DOI 10.4171/JEMS/1633であることを確認した．今回，購読が必要な出版版本文には
アクセスできなかったため，出版版との番号対応は断定しない．
Gomesは改題後のarXiv:2604.04134v2，2026年5月21日を指す．
同論文の実次元$2r+2$における$r$は，本稿の$m=n-1$であり，最大値$2r+1$は
本稿の$2n-1$である．Theorems 1.1, 6.6は左不変複素構造による分類を与え，
葉のファイブレーションをAlbanese写像と同定している．その証明は$J$の変形を用い，
Theorem 4.3で元に戻すための独立した議論を行う．本稿では，変形した複素構造の結論を
元の構造に移すのではなく，$(M,J)$にIstratiの定理を直接適用する．
\end{rem}

\section{定義と横断Hodge理論}\label{sec:defs}
$\omega(X,Y)=g(JX,Y)$とし，実1形式に対して$J\alpha=-\alpha\circ J$と定める．
$g$を正の定数$|\theta|_g^2$倍すると$|\theta|=1$となり，$J$，$\theta$，
およびLevi-Civita接続についての平行性は変わらない．
\begin{equation}\label{eq:vaisman}
U=\theta^\sharp,\quad V=JU,\quad\eta=J\theta,
\quad\wt=-d\eta,\quad \omega=\wt+\theta\wedge\eta
\end{equation}
とおく．$\wt$は閉かつ半正定値で，複素rankは$m$，核は
$T\cF=\langle U,V\rangle_{\R}$である．これが横断K\"ahler構造を定める．
これらの標準的なVaisman恒等式は\cite[Section 2.1]{Istrati}，
\cite[Theorem 2.2]{IO}にも記載されている．

\begin{defn}\label{def:basic}
微分形式$\alpha$がbasicであるとは，$\cF$に接するすべてのベクトル場$X$について
$\iota_X\alpha=\cL_X\alpha=0$となることをいう．これは$U,V$について水平かつ
不変であることと同値である．複体$(\Omega_{\bas}^{\bullet},d_{\bas})$のcohomologyを
$H_{\bas}^{\bullet}(M,\R)$と書く．複素化は横断的な型に分解され，
$d_{\bas}=\partial_{\bas}+\bar\partial_{\bas}$となる．
\begin{align}\label{eq:cohomdef}
H_{\bas}^{p,q}&=\frac{\ker(\bar\partial_{\bas}:\Omega_{\bas}^{p,q}
\to\Omega_{\bas}^{p,q+1})}{\bar\partial_{\bas}\Omega_{\bas}^{p,q-1}},\\
H_{\BC,\bas}^{1,1}(M,\R)&=
\frac{\{\alpha\in\Omega_{\bas}^{1,1}:d\alpha=0,\ \bar\alpha=\alpha\}}
{\{i\partial_{\bas}\bar\partial_{\bas}f:f\in C^\infty_{\bas}(M,\R)\}}
\notag
\end{align}
と定義する．basic形式とbasic関数を通常の形式と関数に置き換えると，
$H_{\BC}^{1,1}(M,\R)$が得られる．
\end{defn}

横断正則接束は$Q^{1,0}=(TM/T\cF)^{1,0}$で，rankは$m$である．
横断正則座標におけるHermitian計量$h_T$のChern接続行列は
$h_T^{-1}\partial_{\bas}h_T$であり，その曲率を$F_T$とする．不変な局所行列は
basic接続として貼り合わさる．これは葉方向のBott部分接続を延長する接続である．
\begin{equation}\label{eq:chern}
c^{\bas}(Q^{1,0})=
\left[\det\left(I+\frac{i}{2\pi}F_T\right)\right]_{\bas},
\quad \rhoT=i\tr F_T=-i\partial_{\bas}\bar\partial_{\bas}\log\det h_T
\end{equation}
と定める．basic Hermitian計量の変更によるChern形式の差はbasic exactである．
$TM$について同じ規約を用いると，$\rhoC=-i\partial\bar\partial\log\det g$，
$c_1^{\BC}(M)=[\rhoC/(2\pi)]_{\BC}$となる．
したがって本稿で扱うのは\emph{接束}のChern類であり，標準束の第一Chern類は
その符号を反転したものである．テンプレートのマクロ$\ddbar$には$1/(2\pi)$が
含まれるが，$i\partial\bar\partial$には含まれない．

\begin{lem}[basicポテンシャル補題]\label{lem:ddbar}
canonical foliationは横断的に向きづけられたRiemannian K\"ahler foliationであり，
葉は極小である．basic $\partial\bar\partial$補題が適用できる．特に，実basic
$d_{\bas}$-exactな$(1,1)$形式$\alpha$は，実basic関数$f$を用いて
$\alpha=i\partial_{\bas}\bar\partial_{\bas}f$と書ける．
\end{lem}
\begin{proof}
$U,V$は可換な実正則Killing場で，長さは1である．$U$の平行性から
$\nabla_UU=\nabla_VU=0$，可換性から$\nabla_UV=0$となる．
また定長Killing場の恒等式から
$\nabla_VV=-\tfrac12\operatorname{grad}|V|^2=0$である．
したがって葉は全測地的で，平均曲率は消える．計量はbundle-likeであり，
$\wt^m$が法束を向きづける．よってbasic Hodge理論とK\"ahler恒等式が使える．
\cite[Sections 3.3--3.4, Proposition 3.5.1]{EK}，\cite[Section 2.2]{IO}を参照．
これらの仮定は，大域的な商orbifoldが存在する前から成立している．

必要な次数2の帰結については，次のHodge理論的な証明を書ける．
$\chi=\theta\wedge\eta$とおくと$d\chi=\theta\wedge\wt$である．
basic exact形式を$\wt^{m-1}\wedge\chi$に掛けて積分すると，Stokesにより0となる．
余分な項は横断次数$2m+1$の形式を含むため消える．したがって
$\Lambda_T\alpha$の平均は0である．basicな関数のLaplacian，すなわち非零の
定数倍を除いて$f\mapsto\Lambda_T(i\partial_{\bas}\bar\partial_{\bas}f)$は，
自己共役で，核は定数関数からなる．コンパクトなbasic Hodge理論により
\[
\Lambda_T(i\partial_{\bas}\bar\partial_{\bas}f)=\Lambda_T\alpha
\]
を解ける．差$\beta$は閉かつprimitiveである．$m\geq2$のとき，横断Hodge starは
$*_T\beta=-\beta\wedge\wt^{m-2}/(m-2)!$なので，$\beta$はbasic coclosedでもある．
basic exactでもあることから，そのノルムの二乗は0となる．$m=1$ではtrace-freeな
$(1,1)$形式自体が0である．これで主張が示された．同じbasic Hodge理論から
$H_{\bas}^{2m}(M,\R)=\R$も従う．これはEl Kacimi-Alaouiの定式化で必要な
homological orientabilityである．
\end{proof}

\begin{defn}\label{def:orb}
canonical foliationが\emph{quasi-regular}であるとは葉がコンパクトであることをいい，
有効Lee群がコンパクトな複素1次元群であることと同値である．その作用が自由なとき
\emph{regular}という．quasi-regularの場合，商$\cB$は複素orbifoldである．
正則principal $E$-orbi-bundleとは，局所的に
$(\widetilde W\times E)/\Gamma\to\widetilde W/\Gamma$と書け，遷移関数が
$E$の平行移動として作用する正則束をいう．isotropyもファイバーに平行移動として
作用する．全空間が多様体である場合，対角作用は自由である．通常のprincipal bundleは，
底が滑らかでorbifold isotropyのない場合にあたる．

pure orbifoldとは因子的なisotropyを持たないorbifoldをいう．余次元2以上には
商特異点が残り得る．分類空間$B\cB$を用いて
$H^k_{\orb}(\cB,\Lambda)=H^k(B\cB,\Lambda)$，
$\pi_k^{\orb}(\cB)=\pi_k(B\cB)$と書く．実orbifold cohomologyはorbifold de Rham
cohomologyおよび台空間の実cohomologyと一致するが，整係数では一致するとは限らない．
orbifold被覆とは，局所的に有限な座標群をその部分群に変更する被覆をいう．
台となる特異空間の\'etale被覆とは限らない．多様体の有限\'etale被覆とは，
全空間の有限不分岐正則被覆を意味する．
\end{defn}

LCK rankは$\theta$の周期群のrankであり，同値に$[\theta]$を含む最小の有理部分空間の
次元である．rank 1はregularityともquasi-regularityとも定義が異なる．
構造定理にrank 1の仮定は用いない．

\section{Bott--Chern消滅への橋渡し}\label{sec:bridge}
\begin{lem}[Ricci形式の比較：IstratiのLemma 4.4]\label{lem:ricci}
正規化したVaisman計量について，$M$上の通常の微分形式として
\begin{equation}\label{eq:ricci}
\rhoC=\rhoT
\end{equation}
が成り立つ．
\end{lem}
\begin{proof}
横断正則submersion $p:W\to T$と，$K_T$の局所正則frame $\sigma_T$を取る．
$\theta^{1,0}=(\theta+i\eta)/2$は一般に正則\emph{ではない}．しかし
\[
\sigma=\theta^{1,0}\wedge p^*\sigma_T
\]
は$K_M|_W$の零点のない正則frameである．実際，
$\bar\partial\theta^{1,0}=-i\wt/2$と水平な最高次形式との外積は0になる．
そのノルムの二乗は$|\sigma|_g^2=\tfrac12p^*|\sigma_T|_{h_T}^2$である．
正則な標準束のframeについて，接束のRicci形式は
$i\partial\bar\partial\log|\sigma|^2$と書ける．定数$1/2$は微分すると消えるので
\eqref{eq:ricci}を得る．これにより符号と，余分な$\wt$の倍数が出ないことも確認できる．
\end{proof}

\begin{prop}[三つの条件の比較]\label{prop:bridge}
任意のコンパクト非K\"ahler Vaisman多様体について
\begin{equation}\label{eq:equiv}
c_1^{\bas}(Q^{1,0})=0
\ \Longleftrightarrow\ c_1^{\BC}(M)=0
\ \Longrightarrow\ c_1(M)=0\text{ in }H^2(M,\R)
\end{equation}
が成り立つ．最後の含意は逆にできない．
\end{prop}
\begin{proof}
自然な比較写像による可換図式
\[
\begin{array}{ccc}
H_{\BC,\bas}^{1,1}(M,\R)&\xrightarrow{j}&H_{\BC}^{1,1}(M,\R)\\
\big\downarrow&&\big\downarrow\\
H_{\bas}^{2}(M,\R)&\longrightarrow&H^{2}(M,\R)
\end{array}
\]
がある．左の写像は補題\ref{lem:ddbar}により単射で，像は実横断$(1,1)$部分である．
$c_1^{\bas}=0$なら，この補題を$\rhoT$に適用して
$\rhoT=i\partial_{\bas}\bar\partial_{\bas}f$を得る．basic関数は$M$上の通常の
大域的滑らかな関数であり，それに作用する制限された微分作用素は
$\partial,\bar\partial$と一致する．補題\ref{lem:ricci}からBott--Chern消滅が従う．

逆向きも，quasi-regularityが分かる前に示せる．コンパクトな等長変換群の中で
$U,V$のflowの閉包を$K$とする．これは正則等長変換からなる．実basic $(1,1)$形式が
$\alpha=i\partial\bar\partial f$と書けるとき，Haar確率測度で$f$を$K$について
平均する．$\alpha$の不変性から，平均$\bar f$は
$i\partial\bar\partial\bar f=\alpha$を満たし，$U,V$で不変なのでbasicである．
これは$j$の単射性を示す．$\rhoT=\rhoC$に適用すると\eqref{eq:equiv}の逆向きが従う．
最後の写像はBott--Chernの情報を忘れる写像である．
\end{proof}

忘れられる情報は次のように明示できる．
\begin{lem}\label{lem:lowdegree}
任意のコンパクト非K\"ahler Vaisman多様体について
\begin{align}\label{eq:lowdegree}
H^1(M,\R)&=H^1_{\bas}(M,\R)\oplus\R[\theta],\\
\ker\bigl(H^2_{\bas}(M,\R)\to H^2(M,\R)\bigr)&=\R[\wt]_{\bas}
\notag
\end{align}
が成り立つ．特に，$M$上では$\wt=-d\eta$であるにもかかわらず，
$[\wt]_{\bas}\ne0$である．
\end{lem}
\begin{proof}
$K$の各元は恒等写像とisotopicなので，$K$による平均は通常のde Rham類を変えない．
不変な1形式は，$\beta$をbasic形式，$u,v$をbasic関数として，
$\beta+u\theta+v\eta$と一意に書ける．その外微分は
$d_{\bas}\beta+du\wedge\theta+dv\wedge\eta-v\wt$である．
これがbasicなら，$U,V$との縮約から$du=dv=0$が従う．
$[\wt]_{\bas}$は0ではない．basicなprimitiveが存在すると，
補題\ref{lem:ddbar}と同じStokesの計算により
$\int_M\wt^m\wedge\chi=0$となり，被積分形式の正値性に矛盾する．
したがって閉じた不変1形式は$v=0$を満たし，第一式を得る．和が直和であることは，
不変なexact 1形式$df$について$f$も平均すれば，$U$との縮約が0になることから分かる．
通常の意味でexactなbasic 2形式のprimitiveに同じ議論を行うと第二式が従う．
これは\cite[Sections 3--5]{IO}とも整合する．
\end{proof}

scalar Hopf多様体$(\C^n\setminus\{0\})/\langle q\rangle$，$0<q<1$の葉空間は
$\mP^{n-1}$であり，$c_1^{\bas}=nH\ne0$である一方，$H^2(M,\R)=0$である．
したがって，これは最後の含意の逆に対する反例であって，本稿の主仮定を満たす例
\emph{ではない}．別の言い方をすれば，$c_1(M)=0$はbasic類が0である代わりに
$[\wt]_{\bas}$の非零倍である可能性を許してしまう．

\section{既知の構造定理の適用}\label{sec:structure}
\begin{proof}[定理\ref{thm:main}の構造に関する部分の証明]
命題\ref{prop:bridge}により，IstratiのDefinition 4.1の条件が正確に満たされる．
コンパクト性，Vaisman条件，Lee形式が非零であることはすでに仮定されている．
同論文のCorollary 5.2は，同じ複素多様体上で，指定された許容Lee類をもつ
Chern--Ricci-flat Vaisman計量を与え，これは定数倍を除いて一意である．
補題\ref{lem:ricci}により，その横断計量はRicci-flatである．

同論文のTheorem 5.3は，あらかじめ商を仮定することなく，$\Autop_0(M)$が
コンパクトな複素1次元群であり，Lee群に一致することを示す．
その普遍複素Lie群は$\C$で，核は余コンパクトなrank 2の格子である．
したがって$E=\C/\Lambda$は楕円曲線である．局所自由な作用の軌道が，
元のfoliationのコンパクトな葉となり，ここで初めて商を取ることができる．
この段階がquasi-regularityを証明し，K\"ahler orbifold商を循環的に仮定することを避ける．

同定理の仕組みを説明すると，canonical計量の一意性によって$\Autop_0(M)$の各元が
等長変換になる．各元がLee類と全体積を保つからである．次にコンパクト自己同型群に
関する結果によって，それをLee群と同定する．残りの部分ではorbifold分解定理を適用し，
definiteなSeifert束を用いて，有限被覆後のpurityを証明する．
pure orbifoldに対して用いる分解は\cite[Theorem 6.4, Definition 6.5]{Campana}である．
必要な仮定はコンパクト性，K\"ahler性，orbifold第一Chern類の消滅である．
最後の仮定は，basic形式を商のorbifold形式と同定し，\eqref{eq:hyp}を用いて得られる．
本稿ではIstratiのpurityの議論を既知定理の一部として用い，最初からpurityを置かない．

商にはdefiniteな束に由来する整な正値類があるため，射影的である．
その有限積被覆の各因子も，ample類を制限すれば射影的になる．
したがって平坦な複素トーラスは，単なる複素トーラスではなくabelian varietyである．
単連結性の意味は$\pi_1^{\orb}(Y_i)=\pi_1^{\orb}(Z_j)=0$である．
holonomyはCampanaの定義と同様，pure orbifoldの滑らかな部分のRicci-flat計量の
holonomyを意味する．本稿は$\SU$型因子を$Y_i$と書くので，IstratiのTheorem 5.3とは
二種類の因子の文字の割当てを入れ替えている．$\SU(2)=\Sp(1)$の重複は，
一度だけsymplectic型に分類してよく，二重に数えない．

最後に，商のこの有限被覆に関してorbifold fiber productを取る．局所的には
$(\widetilde W\times E)/\Gamma'\to(\widetilde W\times E)/\Gamma$，
$\Gamma'\subset\Gamma$から誘導される．元の全空間での対角作用が自由なので，
これは多様体の真の不分岐被覆になる．複素構造とVaisman計量は引き戻される．
必要なら連結成分を選んで$M'$を得る．これで\eqref{eq:BB}，\eqref{eq:dimension}が示された．
regularityや因子の滑らかさは主張していない．またIstratiのCorollary 5.4により
$K_M$は正則にtorsionである．
\end{proof}

横断Calabi--Yau定理の入力は，次のようにも表せる．コンパクトでhomologically
orientableな横断K\"ahler foliationについて，$2\pi c_1^{\bas}$の任意の実basic
$(1,1)$代表元は，固定したbasic K\"ahler類に属する一意な横断K\"ahler計量の
Ricci形式になる\cite[Section 3.5.5]{EK}．横断複素次元1では，これは関数の楕円型
方程式に帰着する．$J$を固定したVaisman構造としての実現には
\cite[Theorem 4.1, Lemma 4.2]{OV}とIstratiのCorollary 5.2を用いる．
任意の共形変換で得られたChern--Ricci-flat Hermitian計量が，なおVaismanであるとは
仮定していない．

\section{Betti数，有限被覆，Albanese写像}\label{sec:betti}
\begin{prop}[transgressionの計算]\label{prop:betti}
\eqref{eq:BB}の被覆について
\begin{equation}\label{eq:betti}
b_1(M')=2a+1
\end{equation}
が成り立つ．
\end{prop}
\begin{proof}
principal $E$-orbi-bundleの接続は，特性類
$c\in H^2_{\orb}(B,\Lambda)$を与える．Vaisman接続の曲率は，
$\Lambda\otimes\R$の一つの非零方向における$\wt$の倍数である．
実際，2次元の垂直接空間で$\theta,\eta$を用いると，
$d\theta=0$，$d\eta=-\wt$である．周期の基底を選ぶ操作は，
この二つの曲率成分に可逆な実線形変換を施すだけである．
したがって二成分が$H^2_{\orb}(B,\R)$内で張る空間のrankは1であって，0ではない．
正値性とコンパクト性により，底上で$[\wt]^m\ne0$だからである．

分類空間に対する実Leray--Serre列の次数1の完全列は
\begin{equation}\label{eq:five}
0\longrightarrow H^1_{\orb}(B,\R)\longrightarrow H^1(M',\R)
\longrightarrow H^1(E,\R)\xrightarrow{d_2}H^2_{\orb}(B,\R)
\end{equation}
である．遷移関数が平行移動なので，$H^1(E)$の局所係数のmonodromyは自明である．
$u,v$を周期基底の双対とすると，$d_2(u),d_2(v)$は，transgressionの共通の符号規約を
除いて二つの特性類に等しい．よって$\rank d_2=1$，$\dim\ker d_2=1$である．
orbifold単連結性から$H^1_{\orb}(\Qfac,\R)=0$となるので，
$b_1^{\orb}(B)=2a$である．\eqref{eq:five}の次元を数えて結論を得る．
\end{proof}

\begin{prop}[virtual第一Betti数]\label{prop:virtual}
$b_1^{\virt}(M)$を，連結な有限正則\'etale被覆$N\to M$にわたる$b_1(N)$の上限とする．
この上限は達成され，$2a+1$に等しい．特に
\begin{equation}\label{eq:gap}
b_1^{\virt}(M)\in\{1,3,\ldots,2n-5\}\cup\{2n-1\}
\end{equation}
である．$n=2$では最初の集合を空集合と解釈する．$2n-3$はvirtual不変量の値として
除外されるが，$b_1(M)$の値として除外されるとは限らない．
\end{prop}
\begin{proof}
任意の有限被覆上で，引き戻した横断Ricci-flat計量の制限holonomyは変わらない．
平均曲率が0なので，basic調和1形式$\alpha$に対するbasic Bochner恒等式から
\[
0=\|\nabla^T\alpha\|_{L^2}^2+
\int\Ric^T(\alpha^\sharp,\alpha^\sharp)\,dv
=\|\nabla^T\alpha\|_{L^2}^2
\]
を得る．したがってbasic調和1形式は平行である．非平坦な既約$\SU$型，$\Sp$型の
因子には不変な余ベクトルが存在しない．よって$M'$の任意の有限被覆上の平行な横断
余ベクトルの空間は次元$2a$以下である．補題\ref{lem:lowdegree}から，通常の第一
Betti数は$2a+1$以下となる．$M$の有限被覆については，$M'$とのfiber productの
連結成分を取る．transferにより実cohomologyの引き戻しは単射なので，同じ上界が従う．
命題\ref{prop:betti}により，この値は$M'$で達成される．

非平坦な因子の複素次元は少なくとも2である．複素1次元のRicci-flat K\"ahler
orbifold因子は平坦であり，$A$に含まれるべきものである．この分解における正次元の
単連結因子にはなれない．したがって$a=m$または$a\leq m-2$であり，
\eqref{eq:gap}を得る．Bochner法の部分は\cite[Proposition 6.4]{Gomes}の仕組みでもある．
本稿では，そこで必要となるfoliationの仮定を明示して満たしている．
\end{proof}

\begin{prop}[Albanese写像の記述]\label{prop:alb}
$\dim\Alb(M)=(b_1(M)-1)/2$，$\Alb(M')\cong A$であり，$M'$のAlbanese写像は
平行移動を除いて$\pr_A\pi$である．$M$のAlbanese写像は全射submersionである．
一般にはcanonical foliationによる商写像とは一致しない．
\end{prop}
\begin{proof}
コンパクトVaisman多様体の正則1形式は閉かつbasicであり，$(1,0)$型のbasic調和形式に
正確に一致する．その周期による商は真の複素トーラスである．
\cite[Sections 2--3, Theorem 3.3]{Gomes}，\cite[Section 3]{IO}を参照されたい．
basic Hodge対称性と補題\ref{lem:lowdegree}から次元の式が従う．
$M'$上では，これらの形式は$A$の正則1形式の引き戻しに正確に一致する．
単連結因子には，Ricci-flat Bochner理論により正則1形式が存在しないからである．
ファイバーと$\Qfac$が連結なので，
$\pi_1(M')\to\pi_1^{\orb}(B)\to\pi_1(A)$は全射である．
したがって周期格子は，単にその有理的な広がりだけでなく，整格子として一致する．
Albanese写像の普遍性から，それは$\pr_A\pi$に一致する．ファイバーは$\Qfac$上の
連結な楕円曲線orbi-bundleである．

$M$自身については，Chern--Ricci-flat Vaisman計量を用いると，正則1形式の基底は
横断的に平行な形式からなる．一点での独立性と平行性から，各点で独立になる．
よってAlbanese写像は各点で標的の次元に等しいrankをもつ．像はsubmersionにより開，
コンパクト性により閉であり，連結なAlbaneseトーラス全体に等しい．
\end{proof}

$\widehat M\to M$が群$F$のGalois被覆なら，
$H^1(M,\R)=H^1(\widehat M,\R)^F$である．
有限群が横断的な平行形式を取り除くことが，小さい$b_1$を生じさせる第二の機構である．
これは$\Qfac$があるために$a$が小さくなる機構と区別しなければならない．
$b_1(M)=2n-1$であることと$M$がKodaira多様体であることの同値は
\cite[Corollary 5.11]{Istrati}にある．regularなAlbanese fibrationと左不変複素構造の
記述は\cite[Theorems 1.1, 6.6]{Gomes}による．
\eqref{eq:hyp}がなければ，次元だけによる上界$2n-1$は成り立たない．
第\ref{sec:examples}節の円周束の構成を種数$g$の曲線に適用すると，任意の$g\geq2$に
対して$b_1=2g+1$をもつVaisman曲面が得られる．

\section{特性類と積分解の限界}\label{sec:twist}
$E=\C/\Lambda$に対して，正則principal elliptic orbi-bundleは
$H^1_{\orb}(B,\cO_B(E))$で分類される．完全列
\begin{equation}\label{eq:exp}
0\longrightarrow\Lambda\longrightarrow\cO_B\longrightarrow
\cO_B(E)\longrightarrow0
\end{equation}
が$c(M')\in H^2_{\orb}(B,\Lambda)$を定める．
したがって特性類はデータの一部にすぎない．一つの特性類を固定したとき，その上の
ファイバーは$H^1(B,\cO_B)$の像，すなわち$H^1_{\orb}(B,\Lambda)$の像による商の
下でtorsorとなる．$E$の格子と正則torsorの類も保持しなければならない．

\begin{prop}[線束モデルと再構成]\label{prop:reconstruct}
\eqref{eq:BB}の積被覆$B$上で，primitiveな格子元$e$の向きと，補完する格子元$f$を
選ぶと
\begin{equation}\label{eq:char}
c(M')=\kappa\otimes e,
\qquad \kappa\in H^2_{\orb}(B,\Z),\quad \kappa\text{はK\"ahler類}
\end{equation}
と書ける．$c_1(L)=\kappa$を満たすampleな正則line orbi-bundle $L$と，
$q\in\C^*$，$|q|<1$が存在して
\begin{equation}\label{eq:line}
M'\cong L^{\times}/\langle q\rangle
\end{equation}
となる．さらに
$L\cong\pr_A^*L_A\otimes\bigotimes_i\pr_{Y_i}^*L_i
\otimes\bigotimes_j\pr_{Z_j}^*L_j'$と書ける．
楕円曲線torsorは，同じ$E,q$を用いた各因子のtorsorの引き戻しのBaer和である．
\end{prop}
\begin{proof}
命題\ref{prop:betti}の実曲率の計算では像は1次元である．その整特性類を整2-cycleに
評価すると，この方向に非零の格子ベクトルが含まれることが分かる．したがって，
この直線は有理的である．primitiveなベクトルとその符号を選ぶと，torsionを除いて
正の類$\kappa$が得られる．積被覆では$H_1^{\orb}(B,\Z)=\Z^{2a}$なので，
普遍係数定理から$H^2_{\orb}(B,\Z)$はtorsion-freeである．よってtorsion項は残らず，
\eqref{eq:char}が整係数で成り立つ．

orbifold指数完全列とK\"ahler $(1,1)$定理により，この第一Chern類をもつ
line orbi-bundle $L_0$が得られる．$e$を1に送って
$\Lambda=\Z+\tau\Z$，$\operatorname{Im}\tau>0$と同定する．
すると$E=\C^*/\langle q\rangle$，$q=e^{2\pi i\tau}$であり，
$L_0^{\times}/\langle q\rangle$の特性類は指定したものに等しい．
\eqref{eq:exp}により，これと$M'$との差は$H^1(B,\cO_B)$の像に入る．
line orbi-bundleの指数完全列は，同じ$H^1(B,\cO_B)$を$\Pic^0_{\orb}(B)$に送る．
$L_0$にその像をtensorすることで差を吸収できる．
これで，位相的な類だけが正則束を決めるとは仮定せずに\eqref{eq:line}を示した．
$\Pic^0$によるねじれは正値性を変えない．これは
\cite[Theorem 2.4]{Istrati}で扱われる設定である．

因子分解について，整K\"unneth公式の次数2では，単連結因子を含む交差項が存在しない．
したがって$c_1(L)$は，各因子への制限の引き戻しの和となる．
$L$をそれらの制限のtensor積で割ったものの第一Chern類は0である．
またK\"ahler Hodge理論と$H^1_{\orb}(\Qfac,\C)=0$から
$H^1(B,\cO_B)=H^1(A,\cO_A)$を得る．
残る$\Pic^0$項は$A$から来るので，$L_A$に吸収できる．

具体的には，因子のtorsorを$B$に引き戻したものを$P_k\to B$と書き，
そのfiber productを，平行移動で作用する
\[
\ker(E^r\xrightarrow{\mathrm{sum}}E)
\]
で割る．局所的なファイバーは$E^r/\ker(\mathrm{sum})=E$である．
得られる遷移関数は各因子の遷移関数の和であり，乗法的な$\C^*$座標では線束の
遷移関数の積となる．これは\eqref{eq:line}の遷移関数に正確に一致する．
よって再構成の主張が示された．
\end{proof}

この被覆を取る前には，definiteな束は正のrank 1の実特性類に加えて整torsionを
もち得る．ここでtorsionが消える理由は，特定の積被覆の$H_1$がtorsion-freeだからで
あって，任意のorbifold上のtorsion類を任意の有限被覆でいつでも消せるからではない．
正の類は有限な底の被覆では消えない．実transfer公式
$p_*p^*=(\deg p)\id$が引き戻しの単射性を与えるからである．
ファイバーのisogenyでも，実係数でこの類を消すことはできない．

\begin{prop}[双正則な積に対する，より強い障害]\label{prop:noproduct}
コンパクト非K\"ahler Vaisman多様体は，両因子が正次元のコンパクト複素多様体である
積$X_1\times X_2$と双正則にはならない．この主張は有限\'etale被覆後も成り立つ．
\end{prop}
\begin{proof}
exactかつ半正定値な形式$\wt$を，次元$r_i$のコンパクトslice $X_i$に制限する．
Stokesにより$\int_{X_i}\wt^{r_i}=0$である．半正定値性から，制限のrankは各点で
落ちている．半正定値Hermitian形式では，ノルムが0のベクトルは形式全体の核に入る．
したがって，積の各点で両方の因子について$TX_i\cap\ker\wt\ne0$となる．
$\ker\wt$は複素直線なので，これは相補的な二つの接空間の両方に含まれる必要があり，
矛盾する．有限\'etale被覆もVaismanなので，同じ議論が使える．
これは\cite[Introduction]{Istrati}で言及される積分解の障害を再現する．
\end{proof}

線束モデルでHermitianノルムを選ぶ．零でないベクトルを単位ベクトルと正の長さに
分けると，$L^{\times}$は滑らかに$S(L)\times\R$と同定される．
$q$の偏角が$\varphi$なら，商は$S(L)$上の角度$\varphi$の円周回転のmapping torusである．
この回転は恒等写像とisotopicなので
\begin{equation}\label{eq:smooth}
M'\simeq_{\mathrm{diff}}S(L)\times S^1
\end{equation}
となる．本稿の状況では，isotropyが非零の線上に自由に作用するので，単位円周束の
全空間は滑らかである．適切なノルムを取るとSasaki幾何の構成になる．
\eqref{eq:smooth}は実多様体の微分同相としての積であり，双正則な積を主張しない．
特に$E\times B$の第一Betti数は$2a+2$で，\eqref{eq:betti}とは異なる．
\emph{BB型の類似は底と横断計量については積分解として，全空間については束と群の
拡大による記述として成立する．}

\section{基本群}\label{sec:pi}
分類空間の意味でのhomotopy完全列は
\begin{equation}\label{eq:pi}
\pi_2^{\orb}(B)\xrightarrow{\langle c,\cdot\rangle}\Lambda
\longrightarrow\pi_1(M')\longrightarrow\Z^{2a}\longrightarrow1
\end{equation}
である．最初の矢印を省いてはいけない．特に，楕円曲線ファイバーの基本群は
全空間の基本群に単射で入るとは限らない．

\begin{prop}\label{prop:pi}
$\Qfac$が一点なら，$\pi_1(M')$は$\Z$と実$(2a+1)$次元Heisenberg群の格子の積であり，
virtually abelianではない．$\dim\Qfac>0$なら，$\pi_1(M)$はvirtually
$\Z^{2a+1}$である．これらは\cite[Corollary 5.10]{Istrati}の二つの場合に対応する．
\end{prop}
\begin{proof}
$B=A$なら$\pi_2(B)=0$である．偏極の基底を選び，$\kappa$をすべての$d_i>0$からなる
単因子形$(d_1,\ldots,d_a)$にする．中心の円周の向きを選ぶと，表示は
\begin{equation}\label{eq:heis}
\langle x_i,y_i,z,t\mid z,t\text{は中心的},\
[x_i,y_j]=z^{d_i\delta_{ij}},\ [x_i,x_j]=[y_i,y_j]=1\rangle
\end{equation}
となる．動径方向の円周が$t$を与える．この表示は$A$の基本平行四辺形を持ち上げる
ことで得られる．円周方向の巻き数は$c_1(L)$の評価に等しい．交代的な交換子pairingは
$\Q$上で非退化なので，どの有限指数部分群も可換ではない．

$\Qfac$が正次元なら，orbifold Hurewiczにより
$\pi_2^{\orb}(\Qfac)\cong H_2^{\orb}(\Qfac,\Z)$である．
K\"ahler類$\kappa$の制限は非零なので，\eqref{eq:pi}の最初の写像の像は，ある$d>0$に
ついて$d\Z e$である．したがって最後の写像の核は$\Z/d\oplus\Z$となる．
線束モデルでは，\eqref{eq:smooth}により$\Z$因子が滑らかに分離する．
円周束の基本群は$\Z^{2a}$の$\Z/d$による中心拡大である．構造群が円周の平行移動
からなるためである．その交換子は有限である．$d\Z^{2a}$の逆像は可換となる．
$d$倍の元の持ち上げ同士の交換子は，元の交換子の$d^2$乗だからである．
この逆像は有限生成可換群なので，rank $2a$のtorsion-freeな有限指数部分群を含む．
動径方向の$\Z$を戻すと，有限指数を除いて$\Z^{2a+1}$が得られる．
$M'$から$M$に移っても，このvirtualな結論は変わらない．
\end{proof}

\section{具体例と境界例}\label{sec:examples}
\begin{eg}[具体例に用いる構成]\label{eg:construction}
$B$を滑らかな射影多様体，$H$をample線束とし，$\kappa=c_1(H)$，$L=H^{-1}$とおく．
$i\partial\bar\partial\log h=2\pi\kappa_h>0$を満たす$L$のHermitian計量$h$を選び，
$L^{\times}$上で$F(v)=|v|_h^2$とおく．このとき
\begin{equation}\label{eq:cone}
\Omega_K=i\partial\bar\partial F>0,\qquad
\Omega_V=F^{-1}\Omega_K,\qquad\theta=-d\log F
\end{equation}
である．実際，水平部分は$2\pi F\kappa_h$，垂直部分は
$iF\partial\log F\wedge\bar\partial\log F$となり，各々の方向に正である．
よって$d\Omega_V=\theta\wedge\Omega_V$となる．対数動径座標において，計量は積
$dt^2+g_{S(L)}$の定数倍なので，$\theta$は平行である．$0<q<1$による乗法は
$\Omega_V,\theta$を保つため，商$M(H,q)=L^{\times}/\langle q\rangle$は
コンパクトVaisman多様体となる．canonicalな葉は
$E=\C^*/\langle q\rangle$ファイバーで，商は$B$である．
ファイバーを反時計回りに向きづけると
\begin{equation}\label{eq:exchar}
c(M(H,q))=(-\kappa,0),\qquad
c_j^{\bas}(M(H,q))=c_j(TB)
\end{equation}
となる．最初の周期生成元の符号を変えると，\eqref{eq:char}の正の規約になる．
反転$v\mapsto v^{-1}$は，$\C^*$作用の向きを逆にして$H$と$H^{-1}$の零切断を除いた
束を同定するので，この符号の選択はample線束のモデルと整合する．
非零全空間が滑らかである正のline orbi-bundleにも，同じ局所的構成を適用できる．
\end{eg}

\begin{eg}[$b_1$が最大：Kodaira多様体]\label{eg:torus}
$B=A$を次元$m$の偏極abelian varietyとし，例\ref{eg:construction}を適用する．
得られる多様体の次元は$n=m+1$で，regularであり，すべての$j>0$について
$c_j^{\bas}=0$，また$b_1=2m+1=2n-1$となる．基本群は，必要なら$z$の符号を変えて
\eqref{eq:heis}で与えられる．Albanese写像は楕円曲線束の射影である．
これらは\cite[Example 7.2, Corollary 5.11]{Istrati}のモデルに正確に一致し，
\cite[Theorem 6.6]{Gomes}とも整合する．標準モデル計量はLCK rank 1をもつが，
これは同じ複素多様体上のすべてのVaisman計量についての主張ではない．
$m=1$ならprimary Kodaira曲面である．
\end{eg}

\begin{eg}[K3を底とし$b_1=1$となる例]\label{eg:k3}
$S\subset\mP^3$を滑らかな四次K3曲面とし，$H=\cO_S(1)$とおく．
随伴公式から$K_S\cong\cO_S$であり，Lefschetz超平面定理により$S$は単連結である．
超平面類$h$はprimitiveである．実際，$h^2=4$であり，$h=2u$と書けたとすると
$u^2=1$となってK3曲面の交点形式が偶であることに反する．
それより大きな可除性は数値的に不可能である．Poincar\'e双対性から交点pairingは
unimodularなので，$h$を$H_2(S,\Z)$に評価する写像は$\Z$へ全射である．
ここで用いたK3のcohomologyと格子の事実については
\cite[Chapters 1 and 14]{Huybrechts}を参照されたい．

$M=M(H,q)$に対し，$n=3$で，canonical foliationは底$S$をもつregularなfoliationであり，
\begin{equation}\label{eq:k3numbers}
c_1^{\bas}=0,\quad \int_S c_2^{\bas}=24,\quad
b_1(M)=b_1^{\virt}(M)=1,\quad \pi_1(M)=\Z
\end{equation}
が成り立つ．$c_2$は$c(TS)=(1+h)^4/(1+4h)$から直接計算でき，次数2の項は$6h^2$，
その積分は24である．円周束の完全列
$\pi_2(S)\to\Z\to\pi_1(S(L))\to0$とprimitivityから
$\pi_1(S(L))=0$が従う．動径方向の円周が，表示した基本群を与える．
特性類の組は$(-h,0)$である．葉空間はK3曲面だがAlbaneseは一点である．
滑らかな多様体としては$M=S(L)\times S^1$であり，双正則な意味では非自明な
コンパクト積分解をもたない．
\end{eg}

\begin{eg}[有限群によって$b_1$が小さくなる例]\label{eg:secondary}
楕円曲線$C$，対称な負の線束$L=\cO_C(-[0])$，対合$\iota:x\mapsto-x$を取る．
これは$\beta^2=1$となる線束の持ち上げ$\beta$をもつ．任意の持ち上げを定数の平方根で
rescaleすればよい．$\beta$不変な負のHermitian計量と$0<q<1$を選び，
\begin{equation}\label{eq:secondary}
P=L^{\times}/\langle q^2\rangle,
\qquad X=L^{\times}/\langle q\beta\rangle
\end{equation}
とおく．$P$に誘導される対合は自由である．固定点があれば，ある$k\in\Z$について
$q\beta(v)=q^{2k}v$となるが，ノルムを取ると$q=q^{2k}$となって矛盾する．
不変計量が降下するので，$P\to X$はVaisman曲面の\'etale二重被覆である．
$H^1(P,\R)=H^1(C,\R)\oplus\R[\theta]$において，対合はそれぞれ$-1$と$+1$で作用する．
したがって
\begin{equation}\label{eq:secondarynumbers}
b_1(X)=1,\qquad b_1(P)=b_1^{\virt}(X)=3
\end{equation}
である．$X$の葉空間は平坦orbifold
$C/\{\pm1\}=\mP^1(2,2,2,2)$なので，$c_1^{\bas}(X)=0$である．
これはquasi-regularだがregularではない．特性類を$C$に引き戻すと$(-[0],0)$となり，
実係数部分は非零のrank 1をもつ．局所的な位数2の固定部分群は，
$\C^*/\langle q^2\rangle$の元$q\beta|_{L_x}$による平行移動としてファイバーに作用する．
これは，引き戻した特性類に加えてisotropyのデータを指定している．

対応する標準的な正方トーラスのモデルでは，基本群は
\[
\pi_1(X)=\langle x,y,z,s\mid z\text{は中心的},\ [x,y]=z,
\ sxs^{-1}=x^{-1},\ sys^{-1}=y^{-1},\ szs^{-1}=z\rangle
\]
である．同値に，生成元$s$が水平方向の反転として作用する
$H_{\Z}\rtimes\Z$である．$x,y,z,s^2$で生成される部分群は
$H_{\Z}\times\Z$で，指数は2である．このモデルは，群演算
$(x,y,z)(x',y',z')=(x+x',y+y',z+z'+xy')$をもつ実Heisenberg群で直接見られる．
水平反転$(x,y,z)\mapsto(-x,-y,z)$は整格子と標準Sasaki構造を保つ．
動径周期の半分だけ平行移動することで商の作用が自由になる．
これにより，誘導される次数1の偏極に対して上記の正則構成が実現される．
これらは\cite[Example 7.3]{Istrati}と同じsecondary Kodaira型の例である．
\end{eg}

\begin{eg}[有限被覆で滑らかにできないorbifold因子]\label{eg:orb}
重み付き次数5の多項式
\[
f=z_0^2z_1+z_1^5+z_2^5+z_3^5,
\qquad(w_0,w_1,w_2,w_3)=(2,1,1,1)
\]
を考える．$W=\{f=0\}\setminus\{0\}$は滑らかである．
四つの偏微分の共通零点が原点だけだからである．重み付き$\C^*$作用の商は
射影K3 orbifold $B\subset\mP(2,1,1,1)$となり，$[1,0,0,0]$に$A_1$ orbifold点を
もつが，因子的なisotropyはない．$5=2+1+1+1$なので，随伴公式からorbifold標準束は
自明である．重み付き乗法による商$W/\langle q\rangle$は，
$c_1^{\bas}=0$を満たす滑らかなコンパクトVaisman三次元多様体である．

これは\cite[Example 7.4]{Istrati}にある明示的なデータである．
複素次元3のこの孤立超曲面特異点のlinkは，超曲面linkの連結性定理
\cite[Chapter 5]{Milnor}によって単連結である．その円周orbi-bundleの完全列から
$\pi_1^{\orb}(B)=0$が従う．したがって$B$には非自明な連結orbifold被覆が一つもなく，
特に有限な滑らかなorbifold被覆は存在しない．全空間は，単連結な原点を除いた錐を
無限巡回群で割ったものなので，$\pi_1=\Z$，$b_1=1$となる．
特性類の組はtautologicalな負のline orbi-bundle $\cO_B(-1)$に対応する
$(-c_1^{\orb}(\cO_B(1)),0)$である．局所群$\Z/2$は，零切断を除いたtautologicalな
線に忠実に作用する．これは構造定理の「orbifold」を「多様体」に置き換えることへの
実際の障害であり，単にその証明を得られていないということではない．
\end{eg}

多様体とその被覆を区別して主要な三例を比較すると，次のようになる．
\begin{center}\small
\begin{tabular}{lllll}\toprule
例 & $n$ & 底 & $b_1$ & $b_1^{\virt}$\\\midrule
Kodaira多様体 & $m+1$ & $A^m$ & $2m+1$ & $2m+1$\\
K3上の構成 & $3$ & $S$ & $1$ & $1$\\
secondary Kodaira & $2$ & $C/\{\pm1\}$ & $1$ & $3$\\\bottomrule
\end{tabular}
\end{center}

\section{本稿で証明する追加の帰結}\label{sec:extra}
\begin{cor}[実現する値]\label{cor:values}
各$n\geq2$について，\eqref{eq:gap}の各値はvirtual第一Betti数として実現する．
\eqref{eq:hyp}のもとで現れる通常の第一Betti数は，正確に$1,3,\ldots,2n-1$である．
\end{cor}
\begin{proof}
virtualな最大値は例\ref{eg:torus}で実現する．$0\leq a\leq m-2$に対して
$k=m-a\geq2$とおき，$Q_k$を$\mP^{k+1}$の滑らかな次数$k+2$の超曲面とする．
随伴公式から$K_{Q_k}\cong\cO$，Lefschetzから単連結性が従う．
そのRicci-flat計量は平坦因子をもたない．そのような因子があれば$H^1$が非零に
なるからである．$A^a\times Q_k$に外部tensor積の偏極を入れ，
例\ref{eg:construction}を適用する．命題\ref{prop:betti}，\ref{prop:virtual}から
$b_1=b_1^{\virt}=2a+1$を得る．

通常の不変量については，楕円曲線の積$A=C^m$と積の対称な偏極を取る．
$0\leq\ell<m$として，$\ell$個の因子で恒等，残りの因子で$-1$となる対合を
持ち上げたものを$\beta$とする．\eqref{eq:secondary}の構成をそのまま適用できる．
動径方向の縮小が，底の固定点の上でも作用の自由性を保証する．
不変な横断1形式の次元は$2\ell$であり，$b_1=2\ell+1$，
virtualな値は$2m+1$となる．$\ell=m$はトーラスを底とするモデル自身である．
すべての場合の奇数性はbasic Hodge対称性と補題\ref{lem:lowdegree}から従い，
上界はすでに証明した．
\end{proof}

\begin{cor}[第二basic Chern類を追加した場合]\label{cor:c2}
\eqref{eq:hyp}のもとで，次は同値である．
\begin{enumerate}
\item 実basic cohomologyにおいて$c_2^{\bas}(Q^{1,0})=0$である．
\item 固定したcanonical foliationと横断複素構造に対し，横断的に平坦な
K\"ahler計量が存在する．
\item $b_1^{\virt}(M)=2n-1$である．
\item Kodaira多様体となる$M$の有限\'etale被覆が存在する．
\end{enumerate}
$n=2$では最初の条件は自動的で，四条件はすべて成り立つ．
\end{cor}
\begin{proof}
商，またはその有限積被覆にRicci-flat計量を入れて考える．
$m\geq2$ではunitary Chern曲率$F$は$\tr F=0$，$\Lambda F=0$を満たす．
skew-Hermitian行列上の$-\tr$が定めるノルムを用いると，primitive形式の
Hodge恒等式から
\begin{equation}\label{eq:c2}
\int_B c_2(TB)\wedge\frac{\omega_B^{m-2}}{(m-2)!}
=\frac1{8\pi^2}\int_B |F|^2\frac{\omega_B^m}{m!}
\end{equation}
を得る．実際，$\tr F=0$なら$c_2=(8\pi^2)^{-1}\tr(F\wedge F)$であり，
$*F=-F\wedge\omega_B^{m-2}/(m-2)!$である．
両式はorbifold座標で適用でき，積分には通常どおり固定部分群の位数の逆数を掛ける．
したがってbasic類が消えると$F=0$となる．横断複素次元1ではRicci-flat性自体が
平坦性を意味する．コンパクト平坦K\"ahler orbifoldは，有限なトーラスorbifold被覆を
もつ．これは先に用いた分解定理の平坦な場合である．その引き戻しがKodaira被覆を与える．
同値に，\eqref{eq:BB}に非平坦因子が存在しない．命題\ref{prop:virtual}から第三の
条件を得る．逆にKodaira被覆上ではbasic Chern類は0である．持ち上げたfoliationの
実transferにより引き戻しは単射なので，下の第二basic類も消える．
同じ分解から，(3)が(4)を導くことも直接分かる．
\end{proof}

これらの帰結は既存の構造定理とHodge理論から本稿で証明したものである．
今回の検索では，文献上新しい結果であることは確認できなかった．新規性は主張しない．

\appendix
\section{依存関係，確認事項，主張しないこと}\label{sec:audit}
論理の流れは，極小な横断K\"ahler foliation，basicポテンシャル補題とRicci比較，
同値\eqref{eq:equiv}，Istratiによる固定$J$での計量の存在とquasi-regularity，
その有限orbifold BB分解，特性類の計算，Betti数・Albanese・基本群の帰結，の順である．

コンパクト性はHodge理論，平均，正値性の積分，自己同型群の議論，Albaneseの像に
用いた．Vaisman条件は極小な葉，正則Killing場，\eqref{eq:vaisman}を与える．
第一basic Chern条件は\eqref{eq:equiv}と横断Ricci-flat性に入る．quasi-regularityは
既知定理で証明した後にのみ用いた．有限被覆は横断的な有限holonomyと指定された
整torsionを除くが，正の曲率類は消さない．homotopy，整cohomology，被覆の議論では
orbifold構造を維持した．射影性は正値性から得られ，代数的な具体例と因子の偏極の
選択に用いた．

中心問題は既知の結果によって解決される．本稿では未解決の段階を意図的に証明済みの
定理として用いてはいない．各因子のtorsor上のすべての有限群作用や，標識を忘れた
全空間同士のすべての同型は分類していない．これらは今後の分類問題であって，
定理\ref{thm:main}の欠けた仮定ではない．すべてのorbifoldが有限被覆で滑らかになる
という主張は用いない．一般のLCK多様体に関する主張も，本稿の証明からは導かない．

監査では\AIstatus{KNOWN}，\AIstatus{DERIVED}，\AIstatus{CANDIDATE}，
\AIstatus{REFUTED}，\AIstatus{GAP}を区別する．
「de Rham消滅ならbasic消滅」，「小さい$b_1$なら非トーラス因子が存在」，
「有限被覆ですべての底が滑らかになる」という偽の含意は，具体例で
\AIstatus{REFUTED}となった．\AIstatus{CANDIDATE}を定理として提示してはいない．
Istratiの出版版の番号を確認できなかったこと，網羅的な新規性調査ではないことは，
書誌上の限界として\texttt{research\_audit\_ja.md}に記録した．
指定された三つのarXiv論文の定理番号は実際の本文で確認した．
Milnorの本における一般の超曲面linkの連結性定理の正確な番号は未確認であり，
標準的な章単位の参照のみを用いた．深い入力については，ここで必要な仮定をすべて
明示した．コンパイルと描画による検証の結果は別のREADMEに記録する．

\clearpage
% Bibliography recovered solely from the corresponding supplied PDF.
\begin{thebibliography}{EKA90}
\bibitem[Cam04]{Campana}
Frédéric Campana. Orbifoldes à première classe de Chern nulle. In \emph{The Fano Conference}, pages 339--351. Università di Torino, 2004. Theorem 6.4 and Definition 6.5 checked in arXiv:math/0402243v1.

\bibitem[EKA90]{EK}
Aziz El Kacimi-Alaoui. Opérateurs transversalement elliptiques sur un feuilletage riemannien et applications. \emph{Compositio Mathematica}, 73(1):57--106, 1990.

\bibitem[Gom26]{Gomes}
Lucas H. S. Gomes. On the classification of Vaisman manifolds with vanishing first basic Chern class and large first Betti number, 2026. arXiv:2604.04134v2, 21 May 2026.

\bibitem[Huy16]{Huybrechts}
Daniel Huybrechts. \emph{Lectures on K3 Surfaces}. Cambridge University Press, 2016. Author's online manuscript, Chapters 1 and 14 consulted.

\bibitem[IO23]{IO}
Nicolina Istrati and Alexandra Otiman. Bott--Chern cohomology of compact Vaisman manifolds. \emph{Transactions of the American Mathematical Society}, 376(6):3919--3936, 2023. Numbered citations refer to the final preprint arXiv:2206.07312v2, 26 September 2022.

\bibitem[Ist25]{Istrati}
Nicolina Istrati. Vaisman manifolds with vanishing first Chern class. \emph{Journal of the European Mathematical Society}, 2025. Published online first, 11 March 2025. Numbered citations in this note refer to arXiv:2304.02582v1, 5 April 2023.

\bibitem[Mil68]{Milnor}
John Milnor. \emph{Singular Points of Complex Hypersurfaces}, volume 61 of \emph{Annals of Mathematics Studies}. Princeton University Press, 1968. Chapter 5; precise theorem number for link connectivity not verified in this run.

\bibitem[OV24]{OV}
Liviu Ornea and Misha Verbitsky. A Calabi--Yau theorem for Vaisman manifolds. \emph{Communications in Analysis and Geometry}, 32(10):2889--2900, 2024. Numbered citations refer to arXiv:2206.08808v2, 24 August 2022.

\end{thebibliography}



\clearpage
\endgroup
\end{document}
